% Modernised reading — fonds Alexandre Grothendieck, shelfmark 151, pages 1-74.
%
% This is an INTERPRETATION, not a transcription. It re-reads the pages in
% today's notation and vocabulary, states the hypotheses the manuscript leaves
% implicit, and drops the critical apparatus so that the mathematics reads
% straight through. Everything the brackets and underlines carried — a doubtful
% reading, a notation he used differently, a missing passage — moves into a
% footnote.
%
% It is held to being mathematically correct as it stands, not merely faithful
% to the page. Where the manuscript is loose or mistaken, the true statement is
% what appears here, and the footnote says what the page has. In any dispute of
% substance the transcription governs: whoever cites Grothendieck must cite the
% batch-NN.fr.tex files.
%
% Scope: the folder was read WHOLE before any of this was written — all four
% batches, pages 1-74, in one pass — and it is written as ONE file for the
% whole shelfmark, which is this project's layout: `151.modern.tex` covers the
% four batches of transcription, exactly as 108.modern.tex, 115.modern.tex and
% 161-3.modern.tex cover theirs.
%
% This matters in one place in particular: pages 2-8 and 16-20 are working
% sheets whose conventions are only fixed on page 21, and the sign conventions
% used for them below are read back from there.
%
% Pass: Opus 5 (claude-opus-5), 2026-08-16 — first pass, unchecked against the
% pages by a human. Merged into a single folder-wide file on 2026-08-22,
% without change to the readings themselves; only the four batch résumés were
% replaced by the single folder résumé below, and their keyword lines united.
\documentclass[11pt,a4paper]{article}
% Le préambule commun à toutes les transcriptions du fonds Grothendieck.
%
% Il tient en une idée : **l'incertitude se compose, elle ne se résout pas.**
% Une transcription qui lisse un mot douteux a détruit la seule chose pour
% laquelle on la faisait — un lecteur doit pouvoir savoir, à chaque endroit,
% ce qui était sur le feuillet et ce qui a été ajouté. Les six macros qui
% suivent sont donc l'appareil critique, et non de la décoration.
%
% Les mêmes commandes sont reconnues par `scripts/render.mjs`, qui produit la
% vue de lecture du site. Une macro ajoutée ici doit l'être aussi là-bas,
% faute de quoi le rendu échouera bruyamment — ce qui est voulu : un rendu
% silencieusement faux est pire qu'un rendu absent.

\ProvidesPackage{grothendieck}[2026/08/08 Transcriptions du fonds Grothendieck]

\usepackage{amsmath,amssymb,amsthm}
\usepackage{mathrsfs}
\usepackage{stmaryrd}
% Les diagrammes commutatifs sont partout dans ce fonds : c'est la langue
% même de la géométrie algébrique de Grothendieck, et une transcription qui
% les rendrait en prose ne transcrirait rien.
\usepackage{tikz-cd}
% Français par défaut — c'est la langue des feuillets — mais l'anglais est
% chargé aussi : la traduction et le résumé sont des documents anglais, et la
% typographie française (espaces avant les deux-points) y serait une faute.
% Ces documents basculent par \selectlanguage{english} en tête de corps.
\usepackage[english,french]{babel}
\usepackage{fontspec}
\usepackage{geometry}
\geometry{a4paper,margin=25mm,marginparwidth=28mm}
\usepackage{marginnote}
\usepackage{xcolor}
% ulem, not soul. `soul`'s \st and \ul are built on a character-by-character
% scan that XeLaTeX's font machinery breaks: the document fails with an
% undefined control sequence rather than degrading. ulem does the same two jobs
% and is Unicode-engine safe. `normalem` stops it hijacking \emph, which the
% transcriptions use for Grothendieck's own emphasis.
\usepackage[normalem]{ulem}
\usepackage{hyperref}
\hypersetup{colorlinks=true,linkcolor=black,urlcolor=black,pdfborder={0 0 0}}

% --- Métadonnées du lot -----------------------------------------------------
% Reprises telles quelles par le script de rendu, qui les affiche en tête de
% la vue de lecture. Elles fixent ce que le fichier prétend couvrir : sans
% elles, une transcription flotte sans point d'attache dans un fonds de seize
% mille feuillets.

\newcommand{\folder}[1]{\gdef\grothFolder{#1}}
\newcommand{\batch}[1]{\gdef\grothBatch{#1}}
\newcommand{\leaves}[2]{\gdef\grothFirst{#1}\gdef\grothLast{#2}}
\newcommand{\foldertitle}[1]{\gdef\grothFoldertitle{#1}}
\gdef\grothFolder{?}\gdef\grothBatch{?}\gdef\grothFirst{?}\gdef\grothLast{?}\gdef\grothFoldertitle{}

% --- Le résumé --------------------------------------------------------------
%
% Il ouvre la lecture modernisée, sous le titre « Résumé », et il est composé
% en petit. Ce n'est pas un ornement : c'est le seul passage du document qui
% ne soit pas des mathématiques, et un lecteur doit pouvoir le reconnaître
% avant de l'avoir lu — puis le sauter s'il connaît déjà le sujet, ou s'y
% arrêter s'il ne le connaît pas. Le filet qui le suit marque la reprise du
% corps.

\newenvironment{resume}
  {\small\setlength{\parskip}{0.45em}}
  {\par\medskip\hrule\medskip}

% --- Mots-clés ---------------------------------------------------------------
%
% En anglais, en clôture du résumé de la lecture modernisée : le vocabulaire
% moderne sous lequel chercher ce que les feuillets construisent. Le manifeste
% du site (`npm run manifest`) les extrait pour étiqueter la cote entière —
% cette ligne est la source unique des tags, il n'y a pas de fichier à côté.
\newcommand{\keywords}[1]{\par\smallskip\noindent
  {\footnotesize\textsc{Keywords} --- \textit{#1}}\par}

% --- L'appareil critique ----------------------------------------------------

% Le numéro de feuillet, en marge. C'est l'ancre qui permet de régler un
% désaccord sur une formule en regardant *un* feuillet plutôt qu'une chemise
% de six cents. Le site s'en sert aussi pour tourner les pages du fac-similé
% quand on fait défiler la transcription.
\newcommand{\leaf}[1]{%
  \marginnote{\footnotesize\color{gray!70}\textsf{#1}}%
  \label{leaf:#1}\ignorespaces}

% Illisible : on ne devine pas. Un mot inventé qui se lit comme les autres est
% le pire résultat possible.
\newcommand{\ill}{\textcolor{red!60!black}{[\,\ldots\,]}}

% Lecture douteuse : le mot est donné, et signalé comme incertain.
\newcommand{\uncertain}[1]{\uline{#1}}

% Ajout de l'éditeur — une lettre manquante, un mot sous-entendu.
\newcommand{\add}[1]{\textcolor{blue!50!black}{[#1]}}

% Biffé par Grothendieck lui-même. Ce qu'il a rayé fait partie du document :
% une rature dit souvent où la pensée a changé de direction.
\newcommand{\struck}[1]{\sout{#1}}

% Note du transcripteur, sur l'état matériel du feuillet.
\newcommand{\note}[1]{\par\smallskip{\small\color{gray!60!black}\textit{[#1]}}\par\smallskip}

% Note marginale de Grothendieck, distincte du corps du texte.
\newcommand{\marginal}[1]{\par\smallskip\noindent\hspace{1em}%
  {\small\color{gray!60!black}#1}\par\smallskip}

% --- Titre ------------------------------------------------------------------

\AtBeginDocument{%
  \begin{center}
    {\large\bfseries Fonds Alexandre Grothendieck --- cote n\textsuperscript{o}~\grothFolder}\\[2pt]
    {\itshape\grothFoldertitle}\\[4pt]
    {\small feuillets \grothFirst--\grothLast\ (lot \grothBatch)}\\[2pt]
    {\footnotesize\color{gray!60!black}Transcription établie d'après le fac-similé numérisé
     par l'Université de Montpellier.\\
     Les lectures incertaines sont soulignées, les illisibles notées [\,\ldots\,],
     les ajouts entre crochets.}
  \end{center}
  \vspace{1em}}

\endinput


\folder{151}
\pages{1}{74}
\dating{[à partir de 1981-à partir de 1990]}
\watermark{Édition de démonstration\\interprétation personnelle de l'œuvre}
\foldertitle{Topos stratifié (1981) : notes manuscrites (s.d.) — lecture modernisée du dossier entier}

\begin{document}

\section*{Résumé}

\begin{resume}
Un espace stratifié est un espace découpé en morceaux : une surface et les
courbes qu'elle porte, une variété algébrique et son lieu singulier, un espace
de modules et son compactifié. Chaque morceau, pris seul, est lisse ; les
singularités n'apparaissent qu'aux jointures. Ce dossier poursuit une question
unique : \emph{si l'on connaît les morceaux, et la manière dont chacun
s'approche des autres, connaît-on l'espace ?}

La réponse naïve — recoller les morceaux le long de leurs intersections — ne
peut pas suffire, et pour une raison qu'on voit sur un dessin. Prenez un cône,
et pour morceaux la pointe et le reste. Leur intersection est vide : la pointe
est fermée, le reste est ouvert, et rien dans ces deux données ne dit que le
cône se referme sur sa pointe. Il faut donc porter, en plus des morceaux, un
enregistrement de la façon dont l'un s'approche de l'autre — un voisinage de la
pointe, et ce qu'il en reste quand on ôte la pointe elle-même. C'est ce que
Grothendieck appelle un \emph{tube}, et tout le vocabulaire du dossier tourne
autour de ce mot : le tube $\mathcal{V}_{i,j}$ d'une strate dans une autre, le
tube épointé $\mathcal{V}^{*}_{i,j}$, et les tubes mixtes qui les interpolent.

Le dossier avance en quatre temps.

Il s'ouvre sur des feuilles de calcul qui ne parlent pas encore de strates. Un
système à trois étages $E_{0}, E_{1}, E_{2}$ muni des permutations des
coordonnées, un texte suivi sur les \emph{hyperrecouvrements} d'un topos indexé
par une petite catégorie quelconque, puis la reprise du même calcul en toutes
dimensions. Ces trois pièces posent la même question sous trois habits :
\emph{quand un système de pièces et de recouvrements multiples détermine-t-il
le tout ?} La condition y prend chaque fois la même forme — les recouvrements
multiples se surjectent sur les produits — et c'est cette forme qui reviendra,
transposée aux strates.

Vient ensuite le texte « Structures stratifiées » lui-même, qui commence par le
cas le plus simple : deux morceaux seulement, un fermé $Y$ dans un espace $X$
et ce qui reste. L'espace est recollé — $X$ est la somme amalgamée du tube et
du complémentaire le long du tube épointé — et de là sortent un théorème de van
Kampen, une suite exacte d'homotopie où apparaissent deux groupes que
Grothendieck baptise, en empruntant à l'arithmétique, \emph{inertie} et
\emph{décomposition}, et une description des faisceaux localement constants sur
chaque morceau comme triplets.

Puis le piège, et c'est le moment le plus instructif du dossier. Tout ce qui
précède ne retient des morceaux que leur groupe fondamental. Grothendieck
calcule ce que ce niveau de données donne pour la sphère
$\mathbb{P}^{1}_{\mathbb{C}}$ — le plan complexe plus un point à l'infini, dont
les morceaux sont un point, un plan et un cercle, tous les trois aussi simples
que possible — et trouve un objet \emph{homotopiquement trivial}, alors que la
sphère ne l'est pas. Il traverse ce moment en une ligne : il avait craint que
le type d'homotopie ne pût pas se récupérer, puis s'est convaincu que la
crainte n'était pas fondée. Ce qui échoue n'est pas le programme, c'est un
invariant trop grossier. La réparation tient en trois pages : le groupe
fondamental est un invariant de dimension~1, et ce qui manquait habite un cran
plus haut — la \emph{gerbe}. Sur la sphère, le compte redevient juste, et l'on
retrouve $H^{2}(S^{2},A) \simeq A$.

Le cadre général s'installe alors. Les morceaux ne sont plus deux mais une
famille indexée par un ensemble ordonné $I$, à laquelle s'attache un objet
combinatoire — les \emph{drapeaux}, points munis d'une chaîne croissante
d'indices — qui s'organise en un espace simplicial dont l'espace est la limite
inductive. C'est vrai et insuffisant, pour la raison que la sphère a déjà
montrée : entre deux morceaux voisins il n'y a qu'un bord, et un bord ne dit
pas comment l'un s'approche de l'autre. D'où le remplacement des fermés par les
tubes, qui sont des voisinages.

Le dernier temps recense le diagramme que ces pièces forment — exactement trois
flèches par couple de strates voisines, et pas une de plus — et énonce, page
68, le \emph{théorème de recollement des tubes} : l'espace se retrouve comme
recollement de ce diagramme. Il est donné en première version et sans
démonstration. La dernière section ouvre l'appareil qui doit la porter, et
c'est de la géométrie de topos : une stratification se transporte par image
inverse, donc le théorème est une question locale ; et les \emph{cribles} — les
parties de $I$ closes vers le bas — répondent aux fermés, ce qui achève un
énoncé dont la deuxième section avait donné la moitié : une stratification est
exactement une application continue de $X$ vers $I$ muni de la topologie où les
ouverts sont les parties croissantes.

Le dossier s'arrête au milieu d'une phrase et d'une construction : la page 74
est la dernière, et la partie qui devait généraliser les tubes mixtes aux
cribles reçoit deux formules et rien d'autre. Cette lecture n'invente pas la
suite.

Deux passages valent d'être lus pour eux-mêmes, parce qu'ils montrent comment
il travaille. À la page 53, il se demande s'il faut exiger que l'intersection
de deux morceaux soit encore un morceau, dessine un cercle coupé en deux arcs
par deux points, et répond que non — la suite le montrera bien. Et à la page
43, revenant sur la théorie d'« intégration » des types d'homotopie qu'il
appelait de ses vœux, il la juge, réflexion faite, relativement triviale et
justiciable de techniques des plus standard. Le jugement est exact : c'est la
théorie des colimites homotopiques, et elle existait déjà.

Où cela mène-t-il ? À ce qu'on appelle aujourd'hui la théorie de l'homotopie
stratifiée : catégorie des chemins sortants, faisceaux constructibles sur un
espace stratifié, colimite homotopique sur le poset des strates. Le programme
que ces pages formulent est celui-là même.

\keywords{hypercovering, Čech nerve, flat functor, Diaconescu theorem, symmetric simplicial object, van Kampen theorem, tubular neighbourhood, homotopy pushout, constructible sheaf, exit-path category, gerbe, poset-stratified space, simplicial space, homotopy colimit, conical stratification, link of a stratum, gluing theorem, sieve, base change, double mapping cylinder, codimension filtration}
\end{resume}
\pagerange{1}{1}
\section*{Le plan de « Structures stratifiées »}

La première page est le sommaire d'un texte en quatre sections, avec la
pagination de l'auteur en marge :

\begin{enumerate}
  \item La situation type — préliminaires d'un programme (p.~1 de l'auteur).
  \item Stratifications globales : le formalisme $\tfrac{1}{2}$ simplicial
    (sans tubes) (p.~10).
  \item Introduction des tubes (p.~15).
  \item Topos canoniques associés à une stratification globale (p.~23).
\end{enumerate}

Ce sommaire est tenu. Le texte lui-même commence à la page 26 du fonds, dans le
lot suivant, et court sans interruption jusqu'à la dernière page du dossier : la
section~1 occupe les pages 26 à 43, la section~2 les pages 44 à 55, la
section~3 les pages 56 à 68, la section~4 les pages 69 à 74, où le dossier
s'arrête au milieu d'une construction.\footnote{La pagination de l'auteur et
celle des archivistes ne coïncident pas : l'auteur numérote ses feuillets, dont
chacun porte deux pages du fonds. Les renvois du sommaire sont à sa
numérotation ; ceux du présent commentaire, comme partout dans cette édition,
sont à celle des archivistes.}

\pagerange{2}{8}
\section*{Le nerf de Čech tronqué : le système $E_{0}, E_{1}, E_{2}$}

Quatre feuilles de travail, numérotées 1 à 4 par l'auteur, calculent dans un
système à trois étages. Le modèle qui les rend toutes transparentes n'est pas
écrit sur la page mais s'y lit dans les formules en coordonnées : partant d'un
ensemble $F$, on pose
\[
  E_{0} = F, \qquad E_{1} = F \times F, \qquad E_{2} = F \times F \times F,
\]
les $p$ oubliant une coordonnée, les $\delta$ en répétant une, et
$\mathfrak{S}_{n+1}$ permutant les facteurs de $E_{n}$ :
\[
  p_{0}(x,y) = x, \quad p_{1}(x,y) = y, \quad \delta_{0}(x) = (x,x),
  \quad \sigma(x,y) = (y,x),
\]
\[
  p_{12}(x,y,z) = (x,y), \quad p_{23}(x,y,z) = (y,z), \quad
  p_{31}(x,y,z) = (z,x),
\]
\[
  \delta_{1}(x,y) = (y,x,x), \quad \delta_{2}(x,y) = (x,y,x), \quad
  \delta_{3}(x,y) = (x,x,y), \quad \rho(x,y,z) = (y,z,x).
\]
C'est le \emph{nerf de Čech} de l'application $F \to \{*\}$, tronqué au
deuxième étage. Ce que les quatre feuilles cherchent est sa présentation : quel
est le plus petit système de données et de relations dont ce modèle soit
l'exemple libre.

\subsection*{La convention sur les transpositions}

Une convention doit être fixée avant tout, car deux relations de ces pages sont
fausses sans elle et vraies avec. Les transpositions $\sigma_{1}, \sigma_{2},
\sigma_{3}$ de $\mathfrak{S}_{3}$ sont indexées par la lettre qu'elles
\emph{fixent} :
\[
  \sigma_{1}(x,y,z) = (x,z,y), \qquad
  \sigma_{2}(x,y,z) = (z,y,x), \qquad
  \sigma_{3}(x,y,z) = (y,x,z).
\]
C'est la convention qu'exigent les relations d'équivariance
\[
  \sigma\, p_{12} = p_{12}\,\sigma_{3}, \qquad
  \sigma\, p_{23} = p_{23}\,\sigma_{1}, \qquad
  \sigma\, p_{31} = p_{31}\,\sigma_{2},
\]
qui sont vraies dans le modèle et sont écrites telles quelles sur la
page.\footnote{La convention n'est explicitée nulle part dans ce lot ; elle
l'est à la page 21, dans le lot suivant, où les trois transpositions sont
données en coordonnées. Nous la lisons donc à rebours, ce que seule une lecture
du dossier entier permet.}

\subsection*{Le système et ses relations}

La donnée minimale est celle de trois ensembles $E_{0}$, $(E_{1},
\mathfrak{S}_{2})$, $(E_{2}, \mathfrak{S}_{3})$ et de quatre applications
$p_{0} \colon E_{1} \to E_{0}$, $\delta_{0} \colon E_{0} \to E_{1}$,
$p_{12} \colon E_{2} \to E_{1}$, $\delta_{1} \colon E_{1} \to E_{2}$ ; tout le
reste s'en déduit par les groupes :
\[
  p_{1} = p_{0}\sigma, \qquad
  p_{23} = p_{12}\rho, \qquad p_{31} = p_{12}\rho^{2}, \qquad
  \delta_{3} = \rho\,\delta_{1}, \qquad \delta_{2} = \rho^{2}\delta_{1}.
\]
Les relations que la page encadre, et qui sont toutes vraies dans le modèle,
sont
\[
  p_{0}\delta_{0} = \mathrm{id}_{E_{0}}, \qquad \sigma\delta_{0} = \delta_{0},
\]
\[
  p_{12}\delta_{1} = \sigma, \qquad
  p_{12}\delta_{3} = \delta_{0}p_{0}, \qquad
  p_{12}\delta_{2} = \mathrm{id}_{E_{1}},
\]
\[
  \sigma_{1}\delta_{1} = \delta_{1}, \qquad
  p_{12}\sigma_{3} = \sigma\,p_{12}.
\]
Les trois du milieu sont les trois valeurs distinctes que prend $p_{12}$ sur
les trois faces $\delta_{1}, \delta_{2}, \delta_{3}$ : une involution, un
idempotent, l'identité. Elles sont, à la lettre, trois des identités
simpliciales.\footnote{La page 2 écrit $\sigma_{3}\delta_{1} = \delta_{1}$ ; la
page 4 écrit $\sigma_{1}\delta_{1} = \delta_{1}$. Sous la convention ci-dessus,
c'est la seconde qui est vraie : $\delta_{1}(x,y) = (y,x,x)$ est fixé par la
transposition des deux dernières lettres, c'est-à-dire par $\sigma_{1}$. Le
stabilisateur de $\delta_{1}(\xi)$ dans $\mathfrak{S}_{3}$ est exactement
$\{1,\sigma_{1}\}$ dès que $\xi$ n'est pas diagonal, et c'est ce fait qui sert
au lot suivant.}

Deux conséquences méritent d'être isolées, parce qu'elles servent plus loin.
La première est que $\delta_{1}\delta_{0}$ est invariante sous $\mathfrak{S}_{3}$
tout entier :
\[
  g\,(\delta_{1}\delta_{0}) = \delta_{1}\delta_{0}
  \qquad (g \in \mathfrak{S}_{3}),
\]
ce qui est immédiat sur le modèle, où $\delta_{1}\delta_{0}(x) = (x,x,x)$. La
seconde est que $p_{0}p_{12}$ est invariante à droite sous $\sigma_{1}$ :
\[
  p_{0}\,p_{12}\,\sigma_{1} = p_{0}\,p_{12} = p_{0}\,p_{13},
\]
$\sigma_{1}$ ne touchant pas la première lettre.\footnote{La page 4 écrit cette
dernière relation avec $\sigma_{3}$ et nomme le membre de droite $p_{1}$ ; les
deux sont fautifs, et de la même façon que la relation
$\sigma_{3}\delta_{1} = \delta_{1}$ de la page 2. Sous la convention fixée,
$p_{0}p_{12}\sigma_{3}$ vaut $p_{1}p_{12}$ et non $p_{0}p_{12}$. La page 6
écrit la relation correcte, avec $\sigma_{1}$.}

\subsection*{Une inégalité qui n'en est pas une}

La page 2 porte, barrée et flanquée de deux points d'interrogation,
\[
  (p_{12}\rho)\,\delta_{1} \neq p_{23}\,(\rho\delta_{1}) \qquad ??
\]
L'inégalité est vraie et n'a rien d'inquiétant. Comme $p_{12}\rho = p_{23}$, le
membre de gauche est $p_{23}\delta_{1} = \delta_{0}p_{0}$ ; le membre de droite
est $p_{23}\delta_{3} = \mathrm{id}_{E_{1}}$. Ce ne sont pas deux
parenthésages du même composé — le second n'est pas
$p_{12}(\rho\delta_{1})$ — et la composition n'est pas en cause. Les deux
valeurs obtenues sont précisément les deux premières lignes de la boîte que la
page écrit juste en dessous.\footnote{Le doute vient des accolades dont l'auteur
surcharge la ligne, qui renomment $\rho\delta_{1}$ en $\delta_{3}$ et
$\rho^{2}\delta_{1}$ en $\delta_{2}$ : un même symbole y sert de nom et de
composé.}

\subsection*{Le cas pointé, et la lecture cocyclique}

Si $E_{0}$ est réduit à un point, $E_{1}$ et $E_{2}$ deviennent des ensembles
pointés par $e_{1} = \delta_{0}(*)$ et $e_{2} = \delta_{1}(e_{1})$, les groupes
respectent les pointages, et $\delta_{0}p_{0}$ devient l'application constante.
Les relations se réduisent à cinq :
\[
  p_{12}\delta_{1} = \sigma, \qquad
  p_{12}\delta_{3} = e, \qquad
  p_{12}\delta_{2} = \mathrm{id}_{E_{1}}, \qquad
  \sigma_{1}\delta_{1} = \delta_{1}, \qquad
  p_{12}\sigma_{3} = \sigma\,p_{12},
\]
où $e$ est l'application constante. La page 8 les récrit multiplicativement,
et c'est le seul endroit du lot où le système soit lu comme une présentation :
\[
  \delta_{0}p_{0}(x) = 1, \qquad
  \sigma(x)\cdot x = 1, \qquad
  p_{12}(z)\,p_{23}(z)\,p_{31}(z) = 1 \quad (z \in E_{2}).
\]
Ce sont, dans l'ordre, la normalisation, l'inversion et la condition de
cocycle. Autrement dit : $E_{1}$ joue le rôle des données de transition d'un
recouvrement, $\sigma$ celui du passage à l'inverse, $\delta_{0}$ celui des
identités, et la troisième relation est la condition de cocycle de Čech sur les
intersections triples.\footnote{La lecture est celle que la page impose par sa
notation multiplicative ; elle n'y est pas nommée. On ajoutera qu'un tel système
avec ses groupes symétriques est, en langue d'aujourd'hui, un \emph{objet
simplicial symétrique} tronqué au deuxième étage — un foncteur sur la catégorie
des ensembles finis non vides à au plus trois éléments — et que le modèle
$E_{n} = F^{n+1}$ en est l'exemple universel.}

Une remarque de marge de la page 8 vaut d'être conservée, car c'est un critère
et non une étape : $\delta_{1}(x)$ est fixé par $\rho$ — donc par
$\mathfrak{S}_{3}$ tout entier — si et seulement si $x = \delta_{0}p_{0}(x)$,
c'est-à-dire si et seulement si $x$ est dégénéré.

\pagerange{10}{14}
\section*{Hyperrecouvrements}

Cinq pages d'un texte suivi, paginées 1 à 5 par l'auteur, demandent ce que doit
être un hyperrecouvrement d'un topos $X$ \emph{indexé par une petite catégorie
$A$ quelconque}, là où la notion reçue est indexée par un objet
simplicial.\footnote{Les hyperrecouvrements de SGA 4, exposé V, et ceux
d'Artin--Mazur sont indexés par des objets simpliciaux ; le texte cite d'ailleurs
Artin--Mazur quelques pages plus loin, à un autre propos. Le passage à une
catégorie d'indices arbitraire est ce que ces pages proposent, et la
reformulation qu'elles atteignent est celle qui rend le passage indolore.}
La donnée est un foncteur
\[
  \varphi \colon A \longrightarrow X, \qquad \alpha \longmapsto U_{\alpha},
\]
et la question est : quelles conditions sur $\varphi$ ?

\subsection*{Les quatre axiomes}

Ils sont posés dans cet ordre.

\begin{itemize}
\item \textbf{Hyp 1.} Les $U_{\alpha}$, $\alpha \in \mathrm{Ob}\,A$, recouvrent
l'objet final $e_{X}$.
\item \textbf{Hyp 1${}'$.} Pour tout $\alpha_{0}$, le foncteur induit
$A/\alpha_{0} \to X/U_{\alpha_{0}}$ satisfait Hyp 1 : les
$U_{\alpha} \to U_{\alpha_{0}}$ indexées par les flèches $\alpha \to
\alpha_{0}$ forment une famille couvrante.
\item \textbf{Hyp 2.} Pour $\alpha, \beta$, la famille des
$U_{\gamma} \to U_{\alpha} \times U_{\beta}$, indexée par les objets de
$A/\alpha\times\beta$ — le produit étant pris dans $\widehat{A}$ —, est
couvrante.
\item \textbf{Hyp 2${}'$.} La version relative : pour tout diagramme
$\alpha_{1} \to \alpha_{0} \leftarrow \alpha_{2}$, la famille des
$U_{\gamma} \to U_{\alpha_{1}} \times_{U_{\alpha_{0}}} U_{\alpha_{2}}$,
indexée par les objets de $A/\alpha_{1}\times_{\alpha_{0}}\alpha_{2}$, est
couvrante.
\end{itemize}

La version à $n$ facteurs de Hyp 2 ne coûte rien : elle se déduit du cas
$n = 1$ par récurrence, puisque la composée de deux familles couvrantes est
couvrante, en écrivant la famille des
$U_{\gamma'} \to U_{\gamma} \times U_{\alpha_{n+1}} \to
(U_{\alpha_{0}} \times \cdots \times U_{\alpha_{n}}) \times
U_{\alpha_{n+1}}$.\footnote{Les deux étiquettes de ce diagramme sont
interverties sur la page relativement à la phrase qui le suit ; la
transcription le signale et ne corrige pas. Rien de mathématique n'en dépend :
la première flèche est celle du cran $n+1$, la seconde celle de l'hypothèse de
récurrence.}

\subsection*{La reformulation, et ce qu'elle dit}

Soit $\varphi_{!} \colon \widehat{A} \to X$ l'unique prolongement de $\varphi$
commutant aux petites colimites — l'extension de Kan à gauche le long de
Yoneda —, et soit
\[
  \varphi^{*} \colon X \longrightarrow \widehat{A}, \qquad
  \varphi^{*}(Y) = \bigl(\alpha \mapsto \mathrm{Hom}_{X}(\varphi\alpha, Y)\bigr)
\]
son adjoint à droite. Le premier est cocontinu, le second continu, et
$\varphi_{!} \dashv \varphi^{*}$.\footnote{L'étoile de $\varphi^{*}$ est ici
celle d'un adjoint à \emph{droite}, à rebours de l'usage : si $\varphi_{!}$
provient d'un morphisme de topos $f$, on a $\varphi_{!} = f^{*}$ et
$\varphi^{*} = f_{*}$. Le manuscrit écrit les deux notations à quelques lignes
l'une de l'autre, page 13 et page 17.}

Alors Hyp 1, Hyp 2, Hyp 2${}'$ sont exactement :
\begin{itemize}
\item a) $\varphi_{!}(e_{\widehat{A}}) \to e_{X}$ est épi ;
\item b) $\varphi_{!}(\alpha\times\beta) \to \varphi_{!}(\alpha) \times
\varphi_{!}(\beta)$ est épi, pour $\alpha,\beta \in \mathrm{Ob}\,A$ ;
\item c) $\varphi_{!}(\alpha\times_{\alpha_{0}}\beta) \to \varphi_{!}(\alpha)
\times_{\varphi_{!}(\alpha_{0})} \varphi_{!}(\beta)$ est épi, pour tout
diagramme $\alpha \to \alpha_{0} \leftarrow \beta$ dans $A$.
\end{itemize}
Et comme tout objet de $\widehat{A}$ est colimite de représentables,
$F = \varinjlim_{A/F}\alpha$, les conditions b) et c) valent alors pour des
$F, G$ quelconques de $\widehat{A}$, ce qui permet de tout ramasser en un
axiome unique :
\[
  \textbf{(Hyp)} \qquad
  \varphi_{!}\Bigl(\varprojlim_{I} F_{i}\Bigr) \longrightarrow
  \varprojlim_{I} \varphi_{!}(F_{i}) \quad \text{est épi},
\]
pour tout diagramme fini $(F_{i})_{i \in I}$ dans $\widehat{A}$ — et il suffit
de le postuler pour $I$ vide, discret à deux objets, et
$({\cdot} \to {\cdot} \leftarrow {\cdot})$, ce qui redonne respectivement
Hyp 1, Hyp 2, Hyp 2${}'$.\footnote{La démonstration que b) et c) sur les
représentables impliquent leur version générale est menée à la page 17, sur des
$\varinjlim$ indexées par $A/F$ et $A/G$ ; elle utilise que ces catégories
d'indices sont filtrantes deux à deux au sens voulu, ce que la page ne dit pas
mais que la construction du produit dans $\widehat{A}$ fournit.}

Ce que dit (Hyp) est net : \emph{$\varphi_{!}$ commute aux limites finies, à
un recouvrement près}. Si l'on y remplace « épi » par « iso », on obtient
l'exactitude à gauche de $\varphi_{!}$ ; le manuscrit relève lui-même ce cas
comme « particulièrement agréable » et note qu'il équivaut à la donnée d'un
morphisme de topos $f \colon X \to \widehat{A}$ avec $\varphi_{!} = f^{*}$.
C'est le théorème de Diaconescu, et ce que le manuscrit appelle le cas agréable
s'appelle aujourd'hui la \emph{platitude} du foncteur
$\varphi$.\footnote{R.~Diaconescu, \emph{Change of base for toposes with
generators}, 1975 : les morphismes de topos $X \to \widehat{A}$ correspondent
aux foncteurs plats $A \to X$, c'est-à-dire à ceux dont l'extension de Kan à
gauche est exacte à gauche. Le résultat est antérieur à ces pages ; rien
n'indique qu'elles s'y réfèrent, et le manuscrit énonce la correspondance sans
la nommer.}

Le texte se termine sur une question, laissée ouverte :

\begin{quote}
Mais cette hypothèse peut-elle être vérifiée, sans que $\varphi_{!}$ soit
exact ? Quand ?
\end{quote}

C'est exactement la question de l'écart entre la platitude et sa forme locale.
L'écart est réel — les deux conditions ne coïncident pas en général, et c'est
pour cela que la seconde a reçu un nom propre, \emph{covering flat} — mais le
dossier n'y revient pas.\footnote{La condition (Hyp) est, mot pour mot, celle
qu'on appelle aujourd'hui platitude relative à la topologie ; elle est ce qui
remplace la platitude dans l'énoncé de Diaconescu lorsque la catégorie
d'indices n'a pas les limites finies. Ces pages en donnent la formulation
correcte sans en tirer d'exemple séparant.}

\pagerange{16}{20}
\section*{Le système $(E_{n}, \partial_{n}, k_{n})$ et son dual}

Une seconde série numérotée, 1 à 3, reprend le calcul en toutes dimensions.
Pour $n \geq 1$, on se donne un $\mathfrak{S}_{n}$-ensemble $E_{n}$ et un
couple
\[
  \partial_{n} \colon E_{n} \longrightarrow E_{n+1}, \qquad
  k_{n} \colon E_{n+1} \longrightarrow E_{n},
\]
les noms venant des applications d'indices sous-jacentes : $\partial_{n}$
répond à l'inclusion $[1,n] \hookrightarrow [1,n+1]$ et $k_{n}$ à sa
rétraction. Les relations sont
\[
  k_{n}\partial_{n} = \mathrm{id}_{E_{n}}, \qquad
  \partial_{n}\sigma = i_{n}(\sigma)\,\partial_{n} \ \ (\sigma \in \mathfrak{S}_{n}),
  \qquad
  \sigma k_{n} = k_{n}\,i'_{n}(\sigma) \ \ (\sigma \in \mathfrak{S}_{n-1}),
  \qquad
  k_{n}\sigma_{n+1} = k_{n},
\]
où $i_{n} \colon \mathfrak{S}_{n} \hookrightarrow \mathfrak{S}_{n+1}$ est
l'inclusion, $i'_{n} = i_{n}i_{n-1}$, et $\sigma_{m}$ la transposition
$(m-1,m)$. La composée dans l'autre sens,
\[
  \rho_{n+1} = \partial_{n}k_{n} \colon E_{n+1} \longrightarrow E_{n+1},
\]
est alors un idempotent, projecteur de $E_{n+1}$ sur l'image de $E_{n}$ : c'est
la seule chose que $k_{n}\partial_{n} = \mathrm{id}$ permette d'affirmer, et
elle suffit.

\subsection*{Le système dual, et la condition de concaténation}

Renversant toutes les flèches, on obtient $E^{*}_{n}$ avec
$d_{n} = k^{*}_{n} \colon E^{*}_{n} \to E^{*}_{n+1}$ et
$p_{n} = \partial^{*}_{n} \colon E^{*}_{n+1} \to E^{*}_{n}$, satisfaisant les
relations transposées. Deux familles d'applications de descente y sont
définies, et il faut les distinguer soigneusement, car c'est leur \emph{couple}
qui porte tout l'énoncé :
\[
  \alpha^{*}_{n,p} = \partial^{*}_{p}\,\partial^{*}_{p+1}\cdots\partial^{*}_{n-1}
  \colon E^{*}_{n} \longrightarrow E^{*}_{p} \qquad (n > p \geq 1),
\]
\[
  \beta^{*}_{n,q} = \partial'^{*}_{q}\,\partial'^{*}_{q+1}\cdots\partial'^{*}_{n-1}
  \colon E^{*}_{n} \longrightarrow E^{*}_{q} \qquad (n \geq q \geq 1),
  \qquad \partial'_{m} = \rho_{m}\,\partial_{m},
\]
$\rho_{m}$ étant la permutation circulaire de $\mathfrak{S}_{m+1}$. Au niveau
des indices, $\alpha_{n,p}$ est l'inclusion du segment initial
$[1,p] \subset [1,n]$ et $\beta_{n,q}$ celle du segment final, $i \mapsto
i + (n-q)$ : $\alpha^{*}$ est donc la restriction au début et $\beta^{*}$ la
restriction à la fin.

La condition que le système doit satisfaire, en plus des relations, est alors
\[
  (**) \qquad
  E^{*}_{p+q} \xrightarrow{\ (\alpha^{*}_{p+q,p},\ \beta^{*}_{p+q,q})\ }
  E^{*}_{p} \times E^{*}_{q} \quad \text{surjectif}, \qquad p,q \geq 1 :
\]
un objet de longueur $p$ et un objet de longueur $q$ étant donnés, il existe
toujours un objet de longueur $p+q$ dont ils sont respectivement le début et la
fin. C'est une condition de concaténation, et c'est le dessin de la page 17 —
deux segments mis bout à bout, avec le point de raccord
marqué.\footnote{Elle a la forme des conditions de remplissage de la théorie
simpliciale, mais on notera que le but est un \emph{produit} et non un produit
fibré : les deux segments d'indices $[1,p]$ et $[p+1,p+q]$ sont disjoints, il
n'y a rien sur quoi les recoller. C'est la même forme que Hyp 2 de la section
précédente, où la famille des $U_{\gamma}$ devait couvrir un produit et non un
produit fibré.}

\subsection*{Le même système écrit en trois termes}

Les pages 18 à 20 redescendent au cas de trois étages et l'écrivent
$X_{1}, X_{2}, X_{3}$, avec $p_{0} \colon X_{2} \to X_{1}$,
$p_{12} \colon X_{3} \to X_{2}$, $\delta_{0} \colon X_{1} \to X_{2}$,
$\delta_{1},\delta_{2},\delta_{3} \colon X_{2} \to X_{3}$ et les groupes
$\mathfrak{S}_{2}$, $\mathfrak{S}_{3}$. C'est le système
$E_{0}, E_{1}, E_{2}$ des pages 2 à 8, décalé d'un
cran.\footnote{Ces $X_{n}$ n'ont rien à voir avec les $X_{i}$ de
« Structures stratifiées », qui sont les parties fermées d'une stratification.
La collision de lettres est réelle et court sur tout le dossier ; nous gardons
ici $E_{0}, E_{1}, E_{2}$ pour le système à trois étages, et réservons $X$ à
l'espace stratifié.}

Deux carrés y sont déclarés cartésiens par l'auteur — il écrit « cart » en leur
centre —, et ils sont ensuite récrits avec $E$ au lieu de $X$ et $\partial$ au
lieu de $\delta$ : le passage d'un système à l'autre est ce que la page cherche.
Le second de ces carrés,
\begin{tikzcd}[column sep=large, row sep=small]
E_{1} & E_{2} \arrow[l, "p_{12}"'] \\
E_{0} \arrow[u, "\delta_{0}"] & E_{1} \arrow[l, "p_{0}"] \arrow[u, "\delta_{1}"']
\end{tikzcd}

dit exactement, une fois $E_{0}$ réduit à un point, que l'image de $E_{1}$ dans
$E_{2}$ est la fibre de $p_{12}$ au-dessus du point marqué. C'est cette
assertion que le lot suivant démontre, à la page 24.

La page 19 pose alors l'hypothèse qui ouvre le lot suivant : $X_{1}$ réduit à un
élément, donc $E_{1}$ un $\mathfrak{S}_{2}$-ensemble pointé et $E_{2}$ un
$\mathfrak{S}_{3}$-ensemble pointé, reliés par
\[
  d \colon E_{1} \rightleftarrows E_{2} \colon p, \qquad p\,d = \mathrm{id}_{E_{1}}.
\]
La page 20 ne porte que trois lignes de calcul en coordonnées, qui reprennent
$d$ et $\rho^{2}d$ ; elles sont reprises et menées à leur terme à la page
21.\footnote{Les deux triplets de la page 20 sont donnés par la transcription
comme des lectures incertaines. Le calcul de la page 21, où les mêmes formules
sont écrites lisiblement, donne $d(x,y) = (x,y,y)$ et
$\rho^{2}d(x,y) = (y,x,y)$ ; c'est cette lecture-là que nous retenons.}

\subsection*{Une page intercalée}

La page 17, écrite au dos de la même lettre administrative que ses voisines,
n'appartient pas à cette série mais au texte « Hyperrecouvrements » : elle
porte le calcul de $\varphi_{!}(F \times G)$ par colimites qui justifie le
passage des représentables aux objets quelconques, et le couple adjoint
$(\varphi_{!}, \varphi^{*})$ écrit en toutes lettres — « $u$ continu, $v$
cocontinu ». Nous l'avons rendue à sa place logique, ci-dessus.
\pagerange{21}{24}
\section*{Le système à trois étages : le cas pointé, achevé}

Les quatre feuilles numérotées 4 à 7 par l'auteur poursuivent, sans coupure, le
calcul des pages 2 à 8 et 16 à 20. Le système est celui de deux ensembles reliés
par un couple section-rétraction,
\[
  d \colon E_{1} \longrightarrow E_{2}, \qquad
  p \colon E_{2} \longrightarrow E_{1}, \qquad p\,d = \mathrm{id}_{E_{1}},
\]
avec $\mathfrak{S}_{2} = \{1,\sigma\}$ opérant sur $E_{1}$, $\mathfrak{S}_{3}$
sur $E_{2}$, un point marqué $e_{1} \in E_{1}$ fixé par $\sigma$, et
$e_{2} = d(e_{1})$.\footnote{L'auteur écrit ici $X_{2}$ et $X_{3}$ pour $E_{1}$
et $E_{2}$, et $X_{1} = \{e_{1}\}$ pour le point ; ces $X$ n'ont aucun rapport
avec les $X_{i}$ de « Structures stratifiées », qui commencent cinq pages plus
loin. Nous gardons la numérotation par degré, $E_{n}$, dans tout le dossier.}
C'est la page 21 qui donne enfin les conventions en coordonnées, valables pour
tout le dossier :
\[
  d(x,y) = (x,y,y), \qquad p(x,y,z) = (x,y), \qquad
  \rho(x,y,z) = (y,z,x),
\]
\[
  \sigma_{1}(x,y,z) = (x,z,y), \qquad \sigma_{2}(x,y,z) = (z,y,x), \qquad
  \sigma_{3}(x,y,z) = (y,x,z),
\]
chaque $\sigma_{i}$ étant donc la transposition qui \emph{fixe} la $i$-ième
lettre.

\subsection*{Une contradiction apparente, et sa levée}

La page encadre trois relations,
\[
  p\,d = \mathrm{id}_{E_{1}}, \qquad p\,\rho\,d = e, \qquad
  p\,\rho^{2}d = \sigma,
\]
où $e$ désigne l'application constante de valeur $e_{1}$ ; et deux relations
d'équivariance, $\sigma_{1}d = d$ et $p\,\sigma_{3} = \sigma\,p$. La seconde
paraît démentie par les coordonnées de la même page, qui donnent
$p\,\rho\,d(x,y) = (y,y)$, lequel n'est pas constant.

Il n'y a pas de contradiction, et la levée est instructive. Dans le modèle en
coordonnées, $E_{n} = F^{n+1}$ pour un ensemble $F$, on a $(y,y) =
\delta_{0}p_{1}(x,y)$ : c'est la diagonale, non une constante. La relation
générale, valable sans hypothèse, est
\[
  p\,\rho\,d = \delta_{0}\,p_{1},
\]
et elle ne devient $p\,\rho\,d = e$ que sous l'hypothèse posée en tête de page,
$E_{0}$ réduit à un point — auquel cas $\delta_{0}p_{1}$ se factorise par un
point et est bien constante. Les deux écritures de la page appartiennent à deux
régimes : les coordonnées au modèle général, les relations encadrées au cas
pointé.\footnote{La transcription signale le désaccord et ne le tranche pas,
comme il se doit. C'est ici, où l'on peut écrire les deux régimes côte à côte,
qu'il se tranche. La même remarque vaut pour la troisième page du lot
précédent, où $p_{12}\rho\delta_{1} = \delta_{0}p_{0}$ est écrit dans le cas
général et « application constante » dans le cas pointé.}

\subsection*{Le sous-système engendré par $d(E_{1})$}

Soit $E'_{2}$ le sous-$\mathfrak{S}_{3}$-ensemble de $E_{2}$ engendré par
$d(E_{1})$. Le stabilisateur de $d(\xi)$ est $\{1,\sigma_{1}\}$ pour tout
$\xi \neq e_{1}$, et $\mathfrak{S}_{3}$ tout entier pour $\xi = e_{1}$ ; d'où
une orbite ponctuelle et des orbites à trois éléments, en bijection avec
$E^{*}_{1} = E_{1} \setminus \{e_{1}\}$ :
\[
  E'_{2} \;\simeq\; \{e_{2}\} \;\amalg\;
  \bigl(\mathfrak{S}_{3} \wedge_{\{1,\sigma_{1}\}} E^{*}_{1}\bigr),
\]
c'est-à-dire l'ensemble induit de $E^{*}_{1}$, muni de l'action triviale de
$\{1,\sigma_{1}\}$, à $\mathfrak{S}_{3}$.\footnote{Le symbole $\wedge$ est celui
de l'auteur pour l'induction ; il l'emploie aussi, quelques pages plus loin, pour
le produit amalgamé de groupes. Ce sont deux opérations différentes sous un même
signe, et nous ne conservons le signe que là où le manuscrit l'écrit.} On pose
alors $E^{*}_{2} = E_{2} \setminus E'_{2}$, sur lequel la restriction
$p^{*}$ de $p$ satisfait encore $p^{*}\sigma_{3} = \sigma\,p^{*}$.

\subsection*{Quand $E_{2}$ recouvre-t-il $E_{1} \times E_{1}$ ?}

La question posée est celle de la surjectivité de
$(p,\ p\rho) \colon E_{2} \to E_{1} \times E_{1}$ — la même condition de
concaténation que le lot précédent rencontrait sous la forme $(**)$. Sur
$E'_{2}$ l'image se calcule exactement :
\[
  (p, p\rho)(e_{2}) = (e_{1},e_{1}), \qquad
  (p, p\rho)\bigl(\rho^{i}d(\xi)\bigr) =
  \begin{cases}
    (\xi,\ e_{1}) & i = 0, \\
    (e_{1},\ \sigma\xi) & i = 1, \\
    (\sigma\xi,\ \xi) & i = 2,
  \end{cases}
\]
pour $\xi \in E^{*}_{1}$. L'image est donc
$\Sigma = \Sigma_{1} \cup \Sigma_{2} \cup \Delta'$, où $\Sigma_{1} = E_{1}
\times \{e_{1}\}$, $\Sigma_{2} = \{e_{1}\} \times E_{1}$ et $\Delta'$ est
l'image de $\delta' = (\sigma, \mathrm{id})$ ; les trois se coupent deux à deux,
et toutes trois ensemble, en le seul point $(e_{1},e_{1})$.

D'où la réponse, complète : $\Sigma = E_{1} \times E_{1}$ si et seulement si
$E_{1}$ a au plus deux éléments. En effet le complémentaire de
$\Sigma_{1} \cup \Sigma_{2}$ est $E^{*}_{1} \times E^{*}_{1}$, et il faut que
$\xi \mapsto (\sigma\xi,\xi)$ y soit surjectif, c'est-à-dire que
$\eta = \sigma\xi$ pour tous $\xi,\eta \in E^{*}_{1}$ ; si $E^{*}_{1}$ avait au
moins deux éléments, deux orbites distinctes sous $\sigma$, ou une orbite à deux
éléments, fourniraient un contre-exemple.

\subsection*{Le carré cartésien}

Il reste alors à examiner le cas $E_{1} = \{e_{1},\xi\}$, où $E_{2} = E'_{2}$ a
quatre éléments. La question est de savoir si le carré
\begin{tikzcd}[column sep=large, row sep=small]
E_{1} & E_{2} \arrow[l, "p"'] \\
E_{0} \arrow[u] & E_{1} \arrow[l] \arrow[u, "\delta_{3}"']
\end{tikzcd}
est cartésien, c'est-à-dire si $p^{-1}(e_{1}) = \delta_{3}(E_{1})$. L'inclusion
$\supset$ est acquise ; l'autre se vérifie sur la classification ci-dessus. Un
élément de $E'_{2}$ distinct de $e_{2}$ s'écrit $\rho^{i}d(\eta)$ avec
$\eta \neq e_{1}$, et $p$ y vaut $\eta$ si $i = 0$, $\sigma\eta$ si $i = 2$ —
tous deux distincts de $e_{1}$. Donc $p(\zeta) = e_{1}$ force $i = 1$, et
$\rho\,d(E_{1}) = \delta_{3}(E_{1})$ : le carré est bien cartésien.

L'identification $\delta_{3} = \rho\,d\sigma$ qui sert ici est écrite sur la
page avec l'exposant de $\rho$ raturé ; le calcul la donne sans ambiguïté,
puisque $d\sigma(x,y) = (y,x,x)$ et $\rho(y,x,x) = (x,x,y) =
\delta_{3}(x,y)$.\footnote{C'est le seul endroit du dossier où un exposant
illisible se laisse restituer par le calcul plutôt que deviné. La transcription
laisse le trou ; nous le comblons, et disons ici par quoi.}

\pagerange{26}{29}
\section*{1. La situation type : un fermé dans un espace}

Le texte annoncé par le sommaire de la page 1 commence ici. La donnée est
\[
  Y \hookrightarrow X, \qquad Y \text{ fermé dans } X,
  \qquad X^{*} = X \setminus Y,
\]
des espaces topologiques, avec l'avertissement qu'il faudra plus tard les
remplacer par des topos.

\subsection*{Trivialité normale locale}

L'hypothèse posée est que la « structure normale » de $X$ le long de $Y$ soit
localement triviale : pour tout $y \in Y$, il existe un voisinage ouvert $X'$ de
$y$ dans $X$, un espace $Z$ pointé par $z$, et un homéomorphisme
$X' \simeq Y' \times Z$ (où $Y' = Y \cap X'$) échangeant l'inclusion
$Y' \hookrightarrow X'$ et $Y' \times \{z\} \hookrightarrow Y' \times Z$ :

\begin{tikzcd}[column sep=large]
Y' \arrow[r, hook] \arrow[d] & X' \arrow[d] \\
Y'\times\{z\} \arrow[r, hook] & Y'\times Z
\end{tikzcd}

L'espace $Z$ est le modèle transverse, et $z$ y est le point singulier : si $Y$
et $X^{*}$ sont des variétés, $Z \setminus \{z\}$ en est une aussi, et $X$ est
« singulier le long de $Y$ ».

\subsection*{Le tube}

On introduit une factorisation canonique de l'inclusion,
\[
  Y \hookrightarrow \mathcal{V}_{Y,X} \hookrightarrow X,
\]
où $\mathcal{V}_{Y,X}$ est le \emph{voisinage tubulaire} — un objet
infinitésimal, dont le manuscrit dit qu'il faudra donner une définition
techniquement souple, et dont il propose, pour $X$ paracompact, le germe de $X$
autour de $Y$ : la limite projective, comme topos, des voisinages ouverts de $Y$
dans $X$.\footnote{Le manuscrit dit que ce germe « est un topos », ce qui est
exact au sens où la limite projective d'un système de topos ouverts en est un ;
c'est aussi ce qui rend la notion utilisable dans le contexte étale, où il n'y a
pas de voisinage ouvert assez petit. La construction reste à cet endroit
programmatique, et le texte y insiste.}

Les deux inclusions ont des comportements opposés, qui sont exactement ce dont
on se servira : $\mathcal{V}_{Y,X} \hookrightarrow X$ a les propriétés d'un
ouvert, et $Y \hookrightarrow \mathcal{V}_{Y,X}$ celles d'une immersion fermée.
Sous les hypothèses de trivialité normale locale — et des conditions de
modération ensemblistes, que le manuscrit invoque sans les
détailler\footnote{Elles sont indispensables : sans elles, le système des
voisinages ouverts peut n'être cofinal dans aucun système raisonnable et la
rétraction n'existe pas. Le texte le sait et l'écrit en marge : l'intuition est
que $\mathcal{V}_{Y,X}$ et $Y$ sont homotopiquement équivalents indépendamment de
la trivialité locale.} — on a une équivalence d'homotopie
\[
  Y \xrightarrow{\ \sim\ } \mathcal{V}_{Y,X}.
\]

Le tube épointé est ce qu'il reste quand on ôte $Y$ :
\[
  \mathcal{V}^{*}_{Y,X} = \mathcal{V}_{Y,X} \cap X^{*}
  = \mathcal{V}_{Y,X} \setminus Y,
\]
et le carré des quatre espaces s'écrit, la flèche du haut étant ce que le
manuscrit appelle joliment le « bouchage de trous le long de $Y$ » :

\begin{tikzcd}[column sep=large]
\mathcal{V}^{*}_{Y,X} \arrow[r, hook] \arrow[d, hook] & \mathcal{V}_{Y,X} \arrow[d, hook] \\
X^{*} \arrow[r, hook] & X
\end{tikzcd}

\subsection*{$X$ comme somme amalgamée}

L'assertion, posée indépendamment de toute hypothèse de trivialité ou de
lissité, est que $X$ s'obtient par recollement de $\mathcal{V}_{Y,X}$ et de
$X^{*}$ le long de $\mathcal{V}^{*}_{Y,X}$ : c'est la somme amalgamée du
diagramme
\[
  X^{*} \longleftarrow \mathcal{V}^{*}_{Y,X} \hookrightarrow \mathcal{V}_{Y,X}.
\]
Au niveau des espaces, l'énoncé est vrai et élémentaire dès que le tube est
remplacé par un vrai voisinage ouvert : si $U$ est un ouvert contenant $Y$,
alors $\{U, X^{*}\}$ est un recouvrement ouvert de $X$, $U \cap X^{*} = U
\setminus Y$, et tout espace est la somme amalgamée de deux ouverts le long de
leur intersection.\footnote{C'est l'hypothèse que le manuscrit ne formule pas et
qui donne à l'énoncé son contenu : la somme amalgamée est prise dans les
espaces, pour le recouvrement ouvert $\{U,X^{*}\}$, et le passage au germe
$\mathcal{V}_{Y,X}$ est un passage à la limite sur $U$. C'est aussi la raison
pour laquelle il tient à ce que $\mathcal{V}_{Y,X} \hookrightarrow X$ « ait les
propriétés d'un ouvert ».} Ce que le manuscrit veut, et qui demande beaucoup
plus, est la version homotopique : que le type d'homotopie de $X$ soit la somme
amalgamée \emph{homotopique} des types d'homotopie des trois autres.

\pagerange{29}{33}
\section*{Van Kampen, et la fibration du tube épointé}

\subsection*{Le carré cocartésien de groupoïdes}

Le premier étage du programme est celui des groupoïdes fondamentaux. Le carré

\begin{tikzcd}[column sep=large]
\Pi_{1}\mathcal{V}^{*}_{Y,X} \arrow[r, "p"] \arrow[d, "i"'] & \Pi_{1}\mathcal{V}_{Y,X} \arrow[d, dashed] \\
\Pi_{1}X^{*} \arrow[r, dashed] & \Pi_{1}X
\end{tikzcd}

doit être cocartésien.\footnote{Les deux flèches qui le ferment sont dessinées
en pointillé sur la page, et les deux termes de droite mis entre parenthèses :
c'est ce que l'énoncé doit fournir, non ce qu'on tient déjà. Nous conservons le
pointillé, qui porte cette distinction.} C'est le théorème de van Kampen sous sa
forme groupoïdale, celle qui n'exige ni connexité ni choix de point base : pour
un recouvrement ouvert par deux pièces, le groupoïde fondamental du tout est la
somme amalgamée des groupoïdes fondamentaux des pièces au-dessus de celui de
l'intersection.\footnote{La forme groupoïdale est due à R.~Brown (1967) ; elle
est antérieure à ces pages, et c'est précisément elle qui rend l'énoncé
utilisable ici, où $\mathcal{V}^{*}_{Y,X}$ n'a aucune raison d'être connexe. Le
manuscrit passe ensuite à la forme pointée, qui suppose $Y$, $X$, $X^{*}$
connexes et un point base choisi — un revêtement universel de
$\mathcal{V}^{*}_{Y,X}$, dit-il, ce qui est en effet la bonne façon de pointer
un groupoïde.}

Sous ces hypothèses de connexité, l'énoncé prend la forme pointée
\[
  \Pi_{1}(X) \;\simeq\;
  \Pi_{1}(X^{*}) \ast_{\Pi_{1}\mathcal{V}^{*}_{Y,X}} \Pi_{1}(\mathcal{V}_{Y,X}),
  \qquad \Pi_{1}(\mathcal{V}_{Y,X}) \simeq \Pi_{1}(Y),
\]
un produit amalgamé de groupes. Tout revient donc à expliciter
$\Pi_{1}(\mathcal{V}^{*}_{Y,X})$ et ses deux images, c'est-à-dire le diagramme

\begin{tikzcd}[column sep=large, row sep=small]
\Pi_{1}(\mathcal{V}^{*}_{Y,X}) \arrow[r] \arrow[d] & \Pi_{1}(Y) \\
\Pi_{1}(X^{*}) &
\end{tikzcd}

\subsection*{Ce que sert la trivialité normale locale}

C'est ici, et seulement ici, que l'hypothèse de trivialité normale locale
travaille : elle doit fournir que l'inclusion
$\mathcal{V}^{*}_{Y,X} \hookrightarrow \mathcal{V}_{Y,X}$ se comporte comme une
\emph{fibration homotopique}, de base $Y \simeq \mathcal{V}_{Y,X}$ et de fibre
le modèle transverse épointé $\mathcal{V}^{*}_{z,Z}$ — le \emph{lien} de la
strate, en langue d'aujourd'hui. Le manuscrit propose de faire de cette
propriété la \emph{définition} de la trivialité normale locale dans les
contextes où la formulation naïve n'a pas de sens, celui des schémas pour la
topologie étale au premier chef ; l'hypothèse de lissité devient alors
techniquement inutile et n'était qu'heuristique.\footnote{La fibre est écrite
$\mathcal{V}^{*}_{z,Z}$ dans la phrase et $\mathcal{V}^{*}_{z,X}$ dans la
formule deux lignes plus bas. C'est $Z$, le modèle transverse, qui est le bon :
$z$ n'est pas un point de $X$. Dans le contexte schématique, ajoute-t-il, il
faudra travailler avec des types d'homotopie profinis et les localiser hors des
caractéristiques résiduelles, comme chez Artin--Mazur.}

D'où la suite exacte d'homotopie de la fibration,
\[
  \cdots \to \Pi_{i}(\mathcal{V}^{*}_{z,Z}) \to
  \Pi_{i}(\mathcal{V}^{*}_{Y,X}) \to \Pi_{i}(Y) \to
  \Pi_{i-1}(\mathcal{V}^{*}_{z,Z}) \to \cdots
\]
qui n'intéresse qu'en basses dimensions. Si la fibre $\mathcal{V}^{*}_{z,Z}$ est
connexe, la flèche $\Pi_{1}(\mathcal{V}^{*}_{Y,X}) \to \Pi_{1}(Y)$ est
surjective et $\Pi_{1}(\mathcal{V}^{*}_{Y,X})$ est une extension de $\Pi_{1}(Y)$
par le quotient de $\Pi_{1}(\mathcal{V}^{*}_{z,Z})$ par l'image de
$\Pi_{2}(Y)$ — donc par $\Pi_{1}(\mathcal{V}^{*}_{z,Z})$ lui-même si
$\Pi_{2}(Y)$ est nul.\footnote{L'hypothèse est écrite sur la page sous la forme
« si $z$ est déconnectant dans $Z$ localement, i.e. si la fibre
$\mathcal{V}^{*}_{z,Z}$ est connexe » ; le mot est signalé comme lecture
incertaine par la transcription, et il dit en français le contraire de la glose
qui le suit — un point qui déconnecte laisse un lien \emph{non} connexe, comme
un hyperplan dans $\mathbb{R}^{n}$. C'est la glose qui porte l'argument, et
c'est elle que nous retenons : la fibre est supposée connexe.}

Lorsque de plus $\Pi_{2}(Y) \to \Pi_{1}(\mathcal{V}^{*}_{z,Z})$ est nulle, on
obtient une extension
\[
  1 \longrightarrow \Pi_{1}(\mathcal{V}^{*}_{z,Z}) \longrightarrow
  \Pi_{1}(\mathcal{V}^{*}_{Y,X}) \longrightarrow \Pi_{1}(Y)
  \longrightarrow 1,
\]
accompagnée de la flèche $\Pi_{1}(\mathcal{V}^{*}_{Y,X}) \to \Pi_{1}(X^{*})$.
Grothendieck baptise le sous-groupe de gauche \emph{inertie} $I$ et le groupe du
milieu \emph{décomposition} $D$ : les noms sont ceux de la théorie des nombres,
et l'analogie est exacte — l'inertie est la monodromie locale autour de la
strate, la décomposition tout ce que le tube épointé
voit.\footnote{Le manuscrit porte les deux noms sous la suite sans dire
explicitement auquel des termes le second se rapporte ; la transcription le
signale. Le rapprochement avec le groupe de décomposition d'une place fixe la
lecture : $D$ est le groupe fondamental du tube épointé, $I$ le noyau de sa
surjection sur celui de la strate.}

Un cas est distingué, et c'est celui dont tout le dossier se sert : celui où les
pièces $Y$, $X^{*}$, $\mathcal{V}^{*}_{z,Z}$ sont des $K(\pi,1)$, c'est-à-dire
ont tous leurs $\Pi_{i}$ nuls pour $i \geq 2$. La suite exacte montre alors
qu'il en est de même de $\mathcal{V}^{*}_{Y,X}$, et bien sûr de
$\mathcal{V}_{Y,X} \simeq Y$.

\pagerange{34}{37}
\section*{Les faisceaux constructibles : la catégorie $\mathfrak{F}$}

Indépendamment de toute hypothèse de nullité ou de trivialité locale, une chose
se reconstruit à partir du seul carré de groupoïdes : la catégorie
$\mathfrak{F}$ des faisceaux $F$ sur $X$ tels que $F|X^{*}$ et $F|Y$ soient
localement constants. Ce sont les faisceaux qu'on appelle aujourd'hui
\emph{constructibles} pour la stratification à deux strates.

Cette catégorie est équivalente à celle des triplets
\[
  \bigl(E_{X^{*}},\ E_{Y,X}\,;\ \varphi\bigr),
\]
où $E_{X^{*}}$ est un système local sur $\Pi_{1}X^{*}$, $E_{Y,X}$ un système
local sur $\Pi_{1}\mathcal{V}_{Y,X} \simeq \Pi_{1}Y$ — c'est-à-dire un
revêtement étale de $Y$ —, et où $\varphi$ est un morphisme de systèmes locaux
sur $\Pi_{1}\mathcal{V}^{*}_{Y,X}$ :
\[
  \varphi \colon p^{*}(E_{Y,X}) \longrightarrow i^{*}(E_{X^{*}}).
\]
Lorsque les trois groupoïdes sont connexes, cela se dit sans groupoïdes :
$E_{X^{*}}$ est un ensemble à opérateurs $\Pi_{1}(X^{*})$, $E_{Y,X}$ un
ensemble à opérateurs $\Pi_{1}(Y)$, et $\varphi$ une application entre eux
compatible aux opérations restreintes à $\Pi_{1}(\mathcal{V}^{*}_{Y,X})$.

Le manuscrit en tire que $\mathfrak{F}$ est un topos de préfaisceaux,
\[
  \mathfrak{F} \;\simeq\; \widehat{\mathcal{C}},
\]
sur une catégorie $\mathcal{C}$ fabriquée en collant les trois groupoïdes : ses
flèches sont celles des trois groupoïdes, plus, d'un objet $a$ de
$A = \Pi_{1}\mathcal{V}^{*}_{Y,X}$ vers un objet $b$ de
$B = \Pi_{1}\mathcal{V}_{Y,X}$ ou $c$ de $C = \Pi_{1}X^{*}$, les flèches
$p^{*}(a) \to b$ de $B$, respectivement $i^{*}(a) \to c$ de $C$. C'est le
\emph{collage} des deux foncteurs — le cographe du profoncteur de spécialisation
— et c'est, en langue d'aujourd'hui, la troncation en dimension~1 de la
\emph{catégorie des chemins sortants} de l'espace
stratifié.\footnote{La page 37 est recouverte de deux tiers de ratures et la
phrase qui définit $\mathcal{C}$ y est reprise deux fois, avec un mot illisible
au point critique. Nous ne prétendons pas restituer la définition mot pour mot ;
ce qui est sûr, et ce dont les deux calculs qui suivent se servent, est que
$\mathcal{C}$ est obtenue en collant $B$ et $C$ le long de $A$, et que
$\mathfrak{F}$ en est les préfaisceaux.}

\subsection*{Le plan complexe}

Le cas typique pour la géométrie algébrique est celui d'une surface topologique
$X$ et d'un point $Y = \{a\}$ : le tube est le disque infinitésimal $D_{a}$, le
tube épointé le disque épointé $D^{*}_{a}$, et le diagramme des $\Pi_{1}$ se
réduit à

\begin{tikzcd}[column sep=large, row sep=small]
\mathbb{Z}_{a,X} \arrow[r] \arrow[d] & \{1\} \\
\Pi_{1}(X^{*}) &
\end{tikzcd}

où $\mathbb{Z}_{a,X} \simeq \mathbb{Z}$ est le module d'orientation de $X$ en
$a$, dont les deux générateurs correspondent aux deux orientations, et où la
flèche verticale est celle qui définit la structure locale en $a$ du groupe
fondamental de $X^{*}$. Le produit d'un nombre arbitraire de copies de ce modèle
donne les modèles-types des stratifications associées aux diviseurs à croisements
normaux.

Pour $X = \mathbb{C}$ et $a = 0$, la flèche $A \to C$ est un isomorphisme
$\Pi_{1}(D^{*}_{0}) \simeq \Pi_{1}(\mathbb{C}^{*})$ : la catégorie
$\mathcal{C}$ se réduit à deux objets, $\mathrm{Aut}(a) = \mathbb{Z}$ et une
seule flèche $a \to b$, de sorte que $b$ y est objet final. Une catégorie à
objet final a un nerf contractile ; donc $\widehat{\mathcal{C}}$, et avec lui
$\mathfrak{F}$, est de type d'homotopie trivial — ce qui est bien le type
d'homotopie de $X = \mathbb{C}$.

\pagerange{37}{40}
\section*{La sphère, ou pourquoi le niveau $\Pi_{1}$ ne suffit pas}

Prenons maintenant le compactifié :
\[
  X = \mathbb{P}^{1}_{\mathbb{C}} \simeq S^{2}, \qquad
  Y = \{\infty\}, \qquad X^{*} = \mathbb{C}.
\]
Le carré de recollement est

\begin{tikzcd}[column sep=large]
D^{*}_{\infty} \arrow[r, hook] \arrow[d] & D_{\infty} \arrow[d, dashed] \\
\mathbb{C} \arrow[r, dashed] & \mathbb{P}^{1}_{\mathbb{C}} \simeq S^{2}
\end{tikzcd}

et sa somme amalgamée est bien $S^{2}$, dont le type d'homotopie n'a rien de
trivial. Le diagramme correspondant sur les $\Pi_{1}$, lui, est

\begin{tikzcd}[column sep=large, row sep=small]
\mathbb{Z} \arrow[r] \arrow[d] & 1 \\
1 &
\end{tikzcd}

et la catégorie $\mathcal{C}$ qui en sort est la catégorie
${\cdot} \to {\cdot}$, à deux objets et une flèche. Elle a un objet final, donc
$\mathfrak{F} \simeq \widehat{\mathcal{C}}$ est \emph{encore homotopiquement
trivial}, alors que $X$ ne l'est pas — et les pièces de construction
$\{\infty\}$, $\mathbb{C}$, $D^{*}_{\infty}$ sont pourtant toutes des
$K(\pi,1)$, les deux premières étant même contractiles. Le morphisme canonique
\[
  X \longrightarrow \mathfrak{F}
\]
n'est donc pas une équivalence d'homotopie.

\subsection*{Ce que la troncation a perdu}

Le diagnostic se dit en une phrase aujourd'hui, et il vaut la peine de le dire,
parce qu'il montre que rien n'était faux dans la construction — seulement trop
grossier. La catégorie des chemins sortants de $\mathbb{P}^{1}$ stratifié par
$\{\infty\}$ a deux objets, et l'espace des chemins allant de $\infty$ vers le
stratum ouvert est le lien, $D^{*}_{\infty} \simeq S^{1}$. L'espace classifiant
de cette catégorie \emph{à homotopie près} est la somme amalgamée homotopique
$\ast \leftarrow S^{1} \to \ast$, c'est-à-dire $S^{2}$ : rien n'est perdu.
Tronquer en dimension~1, c'est remplacer cet espace de chemins $S^{1}$ par son
$\pi_{0}$, un point — d'où la flèche unique $a \to b$ de $\mathcal{C}$, d'où
l'objet final, d'où la contractilité. C'est exactement ce cran-là que le calcul
de la page 39 met en évidence.\footnote{La formulation moderne — faisceaux
constructibles contre représentations de la catégorie des chemins sortants,
et espace classifiant de celle-ci contre type d'homotopie de l'espace — est
postérieure de plus de vingt ans à ces pages (MacPherson, Treumann, Lurie,
Ayala--Francis--Tanaka). Ce que ces trois pages contiennent est l'observation qui
la rend nécessaire : la troncation en dimension~1 ne suffit pas, et la
raison n'est pas que les pièces soient compliquées, puisqu'elles sont ici les
plus simples possibles.}

Le manuscrit tire lui-même la bonne conclusion, et c'est sur cette phrase que le
lot s'achève :

\begin{quote}
J'avais craint un moment, par cet exemple, que le type d'homotopie de $X$ ne
pouvait pas se récupérer à partir de ceux de $X^{*}$, $Y$,
$\mathcal{V}^{*}_{Y,X}$ et du diagramme de recollement. Mais je me suis
convaincu ensuite que cette crainte n'était pas fondée : le type $\mathfrak{F}$
n'est pas suffisant, lui, pour exprimer le type d'homotopie.
\end{quote}

Ce qui échoue est un invariant, non le programme. Reste à trouver l'invariant
qui marche, et l'essai commence aussitôt, sur le cas de $S^{2}$ : récupérer
$H^{2}(X,A)$ en interprétant cette cohomologie comme classifiant les
\emph{gerbes} sur $X$ liées par $A$. La phrase s'interrompt au bas de la page
40 ; elle se termine à la page 41, dans le lot suivant, et le calcul
réussit.\footnote{Le texte s'arrête au milieu d'une phrase, en bas de feuillet.
Ce que la phrase allait dire n'est pas deviné ici : il est lu à la page 41, qui
la complète.}
\pagerange{41}{43}
\section*{Les gerbes, ou ce que le niveau $\Pi_{1}$ laissait échapper}

La phrase interrompue au bas de la page 40 se termine ici, et elle énonce un
principe de recollement d'un cran plus haut que celui des systèmes locaux : la
donnée d'une gerbe $\mathcal{G}_{X}$ sur $X$ équivaut à celle d'une gerbe
$\mathcal{G}_{X^{*}}$ sur $X^{*}$, d'une gerbe $\mathcal{G}_{Y,X}$ sur le tube
$\mathcal{V}_{Y,X}$, et d'une équivalence entre leurs restrictions au tube
épointé :
\[
  \mathcal{G}_{Y,X}\bigl|\,\mathcal{V}^{*}_{Y,X} \;\simeq\;
  \mathcal{G}_{X^{*}}\bigl|\,\mathcal{V}^{*}_{Y,X}.
\]
C'est le même schéma de recollement que pour les faisceaux constructibles, avec
des gerbes à la place des systèmes locaux — et, comme il s'agit d'une
équivalence entre objets d'une $2$-catégorie et non d'une égalité, c'est
précisément là que loge l'information que la troncation en dimension~1 avait
perdue.

Sur un $K(\pi,1)$, une $A$-gerbe s'interprète canoniquement comme une extension
de $\pi$ par $A$.\footnote{Pour $A$ abélien et le lien donné, les $A$-gerbes sur
$K(\pi,1)$ sont classées par $H^{2}(\pi,A)$, qui classe aussi les extensions de
$\pi$ par $A$ ; l'énoncé du manuscrit est celui-là. Il suppose, comme il le dit,
les pièces connexes et de $\pi_{2}$ nul.} Se donner une gerbe sur $X$ revient
donc à se donner deux extensions — l'une de $\pi_{1}(Y)$ par $A$, l'autre de
$\pi_{1}(X^{*})$ par $A$ — et un isomorphisme entre les deux extensions de $D$
par $A$ qu'elles induisent, $D = \pi_{1}(\mathcal{V}^{*}_{Y,X})$ étant le groupe
de décomposition, muni de ses deux flèches vers $\pi_{1}(Y)$ et
$\pi_{1}(X^{*})$.

\subsection*{Le calcul sur la sphère}

Reprenons $X = S^{2}$, $Y = \{\infty\}$, $X^{*} = \mathbb{C}$, où $D \simeq
\mathbb{Z}$ et où les deux groupes fondamentaux $\pi_{1}(Y)$ et
$\pi_{1}(X^{*})$ sont triviaux. Les deux extensions sont alors triviales, et il
ne reste rien à choisir que l'isomorphisme entre les deux extensions triviales
de $\mathbb{Z}$ par $A$ qu'elles induisent — c'est-à-dire un automorphisme de
l'extension triviale, c'est-à-dire un élément de
$\mathrm{Hom}(\mathbb{Z},A) \simeq A$. D'où
\[
  H^{2}(S^{2}, A) \;\simeq\; A,
\]
la valeur correcte.\footnote{L'interprétation de $H^{2}(X,A)$ comme classifiant
les gerbes liées par $A$ est celle de Giraud (1971), à qui le manuscrit ne
renvoie pas ici. Ce que le calcul établit n'est pas la valeur de
$H^{2}(S^{2},A)$, qui est connue, mais que le schéma de recollement la donne :
c'est un test du programme, non un résultat sur la sphère.}

Ce succès vaut confirmation de l'intuition d'ensemble : il y a une théorie des
$\varinjlim$ de types d'homotopie, et le type d'homotopie d'une limite inductive
de topos est la limite inductive des types d'homotopie. Pour des espaces
stratifiés dont les pièces sont des $K(\pi,1)$ — l'exemple visé étant les espaces
de modules de Teichmüller, leurs compactifiés et leurs voisinages tubulaires —
le type d'homotopie doit s'exprimer en termes d'un système inductif convenable,
le plus souvent fini, des groupoïdes fondamentaux des pièces. C'est le programme
que sert tout le reste du dossier.\footnote{« Cette théorie d'intégration des
types d'homotopie […] est en fait relativement triviale, et justiciable sans
doute de techniques d'explicitation des plus standard », écrit-il, en ajoutant
qu'il se réserve d'y revenir. Le jugement est exact et l'outil existait : les
colimites homotopiques de Bousfield--Kan (1972), et le théorème de Thomason
(1979) identifiant la colimite homotopique d'un diagramme de catégories à la
construction de Grothendieck, sont l'un et l'autre antérieurs à ces pages.}

\pagerange{44}{49}
\section*{2. Stratifications globales : le formalisme $\tfrac{1}{2}$ simplicial, sans tubes}

On quitte le cas de deux morceaux. Le cadre est un espace topologique $X$ — plus
tard, un topos quelconque — et une famille de sous-espaces indexée par un
ensemble ordonné $I$ :
\[
  (X_{i})_{i \in I}, \qquad X_{i} \subset X,
\]
soumise à deux conditions :

\begin{itemize}
\item a) les $X_{i}$ sont fermés et la famille est localement finie ;
\item b) $i \leq j \implies X_{i} \subset X_{j}$.
\end{itemize}

Aucune hypothèse de trivialité locale n'est faite, et les conditions de connexité
et de locale connexité qu'exigera la faisceautisation sur $X$ sont
délibérément différées.

\subsection*{Drapeaux}

On pose
\[
  X_{\Delta_{0}} = \coprod_{i \in I} X_{i}
  = \bigl\{ (x,i) \in X \times I \ \big|\ x \in X_{i} \bigr\},
\]
et la condition a) dit exactement que le morphisme canonique
$X_{\Delta_{0}} \to X$ est \emph{fini} — propre, séparé, à fibres finies ; c'est
aussi une immersion locale. On introduit ensuite la partie fermée
\[
  X_{\Delta_{1}} = \coprod_{i \leq j} \Gamma_{i,j}
  \;\subset\; X_{\Delta_{0}} \times X_{\Delta_{0}},
  \qquad \Gamma_{i,j} \subset X_{i} \times X_{j}
  \text{ le graphe de } X_{i} \hookrightarrow X_{j},
\]
qui n'est autre que le graphe d'une relation d'ordre sur $X_{\Delta_{0}}$ :
\[
  (x,i) \leq (y,j) \iff x = y \text{ et } i \leq j,
\]
soit l'ordre induit par l'ordre produit de $X \times I$, $X$ étant muni de
l'ordre discret. Le couple $(X_{\Delta_{0}}, X_{\Delta_{1}})$ est donc un
\emph{objet ordonné de la catégorie des espaces topologiques} : une catégorie
interne, d'objet des objets $X_{\Delta_{0}}$, d'objet des flèches
$X_{\Delta_{1}}$, de source $\sigma$, de but $b$, et de composition
\[
  X_{\Delta_{2}} = (X_{\Delta_{1}},b) \times_{X_{\Delta_{0}}}
  (X_{\Delta_{1}},\sigma) \longrightarrow X_{\Delta_{1}},
  \qquad (x,y,z) \mapsto (x,z).
\]

Les deux projections ne jouent pas le même rôle, et la dissymétrie est le fait
technique central de la section :

\begin{itemize}
\item $\sigma \colon X_{\Delta_{1}} \to X_{\Delta_{0}}$ est un revêtement étale,
et même constant au-dessus de chaque $X_{i}$, de fibre
$I_{i} = \{ j \in I \mid j \geq i \}$ — donc fini si les $I_{i}$ le sont ;
\item $b \colon X_{\Delta_{1}} \to X_{\Delta_{0}}$ n'est qu'une immersion locale.
\end{itemize}

En degré quelconque, les produits fibrés itérés donnent les drapeaux :
\[
  \begin{aligned}
    X_{\Delta_{r}}
      &= \bigl\{ (x_{0},\ldots,x_{r}) \in X_{\Delta_{0}}^{\,r+1}
        \ \big|\ x_{0} \leq \cdots \leq x_{r} \bigr\} \\
      &\simeq \bigl\{ (x, i_{0},\ldots,i_{r}) \ \big|\
        x \in X_{i_{0}},\ i_{0} \leq \cdots \leq i_{r} \bigr\} \\
      &\simeq \coprod_{i_{*} \in I(\Delta_{r})} X_{i_{0}},
  \end{aligned}
\]
où $I(\Delta_{r}) = \mathrm{Hom}_{\text{ens. ord.}}(\Delta_{r}, I)$ est
l'ensemble des drapeaux de longueur $r$ de $I$. Autrement dit :
$X_{\Delta_{*}}$ est le \emph{nerf} de l'objet ordonné, et il se plonge dans le
produit $X \times I_{\Delta_{*}}$, dont il hérite sa structure simpliciale :
une application croissante $\alpha \colon \Delta_{r'} \to \Delta_{r}$ induit
\[
  \alpha^{*} \colon X_{\Delta_{r}} \longrightarrow X_{\Delta_{r'}},
  \qquad \alpha^{*}(x,i_{*}) = (x,\ i_{*} \circ \alpha).
\]

Deux remarques du manuscrit méritent d'être conservées. D'abord, tout se
généralise en remplaçant l'ensemble ordonné $I$ par une catégorie
$\mathcal{C}$ et la famille par un foncteur de $\mathcal{C}$ vers les parties
fermées de $X$. Ensuite, si l'on part d'une famille sans ordre, on peut toujours
poser $i \leq j \iff X_{i} \subseteq X_{j}$ ; mais la présentation adoptée
n'exige ni que $i \mapsto X_{i}$ soit injective, ni que ce soit un plongement
d'ensembles ordonnés — la réciproque de b) n'est pas requise. C'est une liberté
délibérée, et la section~4 dira pourquoi : elle est ce qui permet de prendre des
images inverses sans avoir à jeter les indices devenus vides.

\subsection*{Ce qui est étale, ce qui ne l'est pas}

La dissymétrie entre $\sigma$ et $b$ se propage à tous les degrés. Si
$\alpha \colon \Delta_{r'} \to \Delta_{r}$ vérifie $\alpha(0) = 0$, alors
$\alpha^{*}$ est un morphisme de revêtements étales au-dessus de $X$ ; en
particulier $X_{\Delta_{r}} \to X_{\Delta_{0}}$, $(x,i_{0},\ldots,i_{r}) \mapsto
(x,i_{0})$, en est un. Plus précisément, au-dessus du composant $X_{i'_{*}}$ de
$X_{\Delta_{r'}}$, le revêtement est constant de fibre l'ensemble des drapeaux
de type $\Delta_{r}$ qui s'envoient sur $i'_{*}$ :
\[
  \alpha^{*-1}\bigl(X_{i'_{*}}\bigr) \;\simeq\;
  X_{i'_{*}} \times \bigl\{ i_{*} \in I(\Delta_{r})
    \ \big|\ i_{*} \circ \alpha = i'_{*} \bigr\}.
\]
Sans la condition $\alpha(0) = 0$, on ne peut affirmer que des immersions
locales — et de même $X_{\Delta_{r}} \to X$ n'est jamais qu'une immersion
locale.\footnote{La raison de la condition $\alpha(0) = 0$ est que le composant
d'un drapeau ne dépend que de son \emph{premier} indice, $X_{i_{*}} = X_{i_{0}}$ ;
une application croissante qui préserve le $0$ ne touche donc pas à l'espace,
seulement à la combinatoire, et c'est ce qui rend $\alpha^{*}$ étale.}

Le manuscrit note enfin qu'on aurait intérêt à éliminer les redondances en se
bornant aux drapeaux \emph{non dégénérés}, c'est-à-dire aux suites strictement
croissantes, ce qui oblige à ne garder que les $\alpha$ strictement croissantes.
C'est la variante semi-simpliciale, et c'est ce que l'auteur appelle
« $\tfrac{1}{2}$ simplicial ».\footnote{Le manuscrit écrit tantôt
« $\tfrac{1}{2}$ simplicial », tantôt « semi-simplicial », pour le même objet.
L'usage d'aujourd'hui distingue : \emph{simplicial} pour le système avec toutes
les applications croissantes, \emph{semi-simplicial} pour la variante à
applications strictement croissantes, qui est celle de son NB de la page 50. Les
deux sont ici en présence, et le texte passe de l'une à l'autre sans le dire.}

\pagerange{50}{53}
\section*{Reconstituer $X$}

La question posée est celle de savoir si le diagramme

\begin{tikzcd}[column sep=large, row sep=small]
& X_{\Delta_{1}} \arrow[dl, "\sigma"'] \arrow[dr, "b"] & \\
X_{\Delta_{0}} \arrow[dr, dashed] & & X_{\Delta_{0}} \arrow[dl, dashed] \\
& X &
\end{tikzcd}

est cocartésien, c'est-à-dire si $X \simeq X_{\Delta_{0}}
\amalg_{X_{\Delta_{1}}} X_{\Delta_{0}}$. Il y faut deux conditions de plus :

\begin{itemize}
\item c) $X = \bigcup_{i} X_{i}$ ;
\item d) pour tous $i,j$, \ $X_{i} \cap X_{j} = \bigcup_{k \leq i,\ k \leq j} X_{k}$.
\end{itemize}

\textbf{Proposition.} Sous a), b), c), d), l'application canonique
\[
  \varinjlim X_{\Delta_{*}} \xrightarrow{\ \sim\ } X
\]
est un homéomorphisme — et l'énoncé vaut pour les deux variantes, drapeaux
quelconques ou drapeaux non dégénérés.

Le corollaire donne la réciproque, et c'est lui qui fait de la construction une
équivalence de données. Soit $(X_{\Delta_{0}}, X_{\Delta_{1}})$ un espace
topologique ordonné tel que $X_{\Delta_{1}}$ soit fermé dans le produit et que,
pour tous $i,j$, la trace $\Gamma_{i,j} = X_{\Delta_{1}} \cap (X_{i} \times
X_{j})$ soit vide ou le graphe d'une immersion fermée $X_{i} \hookrightarrow
X_{j}$. Alors $\sigma$ est un revêtement étale constant, $b$ une immersion
locale, la relation $i \leq j \iff \Gamma_{i,j} \neq \emptyset$ est un ordre sur
$I$, et les $X_{i}$ forment un système inductif d'espaces topologiques à
morphismes de transition des immersions fermées. Si de plus $I$ est fini et
stable par bornes inférieures,
\[
  X = \varinjlim_{i \in I} X_{i} \;\simeq\;
  X_{\Delta_{0}} / \bigl(X_{\Delta_{1}} \rightrightarrows\bigr),
\]
les $X_{i} \hookrightarrow X$ étant des immersions fermées.

\subsection*{Une limite inductive stricte, et pourquoi elle ne suffira pas}

Cette proposition est vraie et n'est pas ce qu'il faudra. La limite inductive
qu'elle exhibe est une limite d'\emph{espaces} : le quotient de
$\coprod_{i} X_{i}$ par la relation d'équivalence engendrée par les inclusions,
c'est-à-dire $\varinjlim_{i \in I} X_{i}$. Elle redonne $X$ comme espace, ce qui
est exact ; mais elle ne dit rien de son type d'homotopie, parce que les
$X_{i}$ sont \emph{fermés} les uns dans les autres. Entre deux morceaux voisins,
la seule chose que ce diagramme enregistre est un bord — et un bord ne dit pas
comment l'un s'approche de l'autre, comme la sphère l'a déjà montré à la section
précédente.

C'est exactement pour cela que la section~3 remplacera les fermés $X_{i}$ par
des \emph{tubes}, qui sont des voisinages. La limite inductive stricte restera
la même ; la limite inductive homotopique, elle, deviendra la bonne.\footnote{La
distinction n'est pas faite sur la page, et le manuscrit n'a pas de mot pour
elle : il écrit $\varinjlim$ dans les deux cas, et parle de « somme amalgamée des
types d'homotopie » quand il veut le second. La théorie qui les sépare — les
colimites homotopiques — lui est contemporaine ; il la nomme lui-même, page 43,
comme la chose qui lui manque et qu'il tient pour standard.}

\subsection*{Faut-il que les intersections soient des morceaux ?}

La condition d) demande que $X_{i} \cap X_{j}$ soit \emph{réunion} des $X_{k}$
en dessous. Faut-il exiger davantage — que la borne inférieure $i \wedge j$
existe et que $X_{i \wedge j} = X_{i} \cap X_{j}$ ?

Le manuscrit répond non, et son contre-exemple est un dessin. Prenez le cercle,
coupé en deux arcs $A$ et $B$ par deux points $a$ et $b$ : quatre strates,
$A \cap B = \{a\} \cup \{b\}$, deux strates connexes dont l'intersection ne l'est
pas. Prenez-en le cône, de sommet $o$ : une stratification locale en $o$ à cinq
strates. « Faut-il l'éliminer pour autant ? Il ne semble pas. La suite le
montrera bien ! » — et il a raison, car ce sont précisément les stratifications
de ce type qui se présentent dans la nature.

\pagerange{53}{55}
\section*{Les strates}

Les $X_{i}$ ne sont pas les morceaux : ils sont les fermés dont les morceaux sont
les différences. Pour tout $x \in X$, l'ensemble $I_{x} = \{ i \in I \mid x \in
X_{i}\}$ est fini par a), et filtrant décroissant par d) ; il a donc un plus
petit élément
\[
  i(x) = \text{le plus petit } i \in I \text{ tel que } x \in X_{i}.
\]
Et pour tout $i$ on pose
\[
  \dot{X}_{i} = \bigcup_{j < i} X_{j} \quad (\text{fermé de } X_{i}),
  \qquad
  X_{i}^{*} = X_{i} \smallsetminus \dot{X}_{i} \quad (\text{ouvert de } X_{i}).
\]

\textbf{Proposition.} Les $X_{i}^{*}$ sont deux à deux disjoints, leur réunion
est $X$, et $i(x)$ est l'unique indice $i$ tel que $x \in X_{i}^{*}$.

\emph{Démonstration.} Soient $i \neq j$ et $x \in X_{i}^{*} \cap X_{j}^{*}$.
Alors $x \in X_{i} \cap X_{j}$, donc par d) il existe $k \leq i,j$ avec
$x \in X_{k}$. Si $k = i$, alors $i \leq j$, donc $i < j$, et $x \in \dot{X}_{j}$,
ce qui contredit $x \in X_{j}^{*}$. Donc $k < i$, d'où $X_{k} \subset \dot{X}_{i}$
et $x \in \dot{X}_{i}$, ce qui contredit $x \in X_{i}^{*}$. La réunion est $X$ et
l'unicité est immédiate par définition de $i(x)$.\footnote{Le signe entre
$\dot{X}_{i}$ et $X_{k}$ est tracé d'un seul trait oblique sur la page et la
transcription hésite ; c'est l'inclusion $X_{k} \subset \dot{X}_{i}$ que
l'argument demande, et c'est elle qui est écrite ici.}

La conséquence à retenir, et que le manuscrit n'énonce pas sous cette forme,
est que $x \mapsto i(x)$ est une application \emph{continue} de $X$ vers $I$
muni de la topologie d'Alexandrov : les parties ouvertes de $I$ y sont les
parties croissantes, et l'image inverse d'une partie décroissante
$\{ j \mid j \leq i \}$ est précisément $X_{i}$, qui est fermé. C'est la
définition moderne d'un espace stratifié par un ensemble ordonné, et elle est ici
entièrement contenue dans les conditions a) à d).\footnote{Cette formulation —
une stratification est un morphisme continu vers un poset alexandrovien — est
celle qui a cours depuis Lurie ; elle est postérieure de trente ans à ces pages.
Ce qui est ici n'est pas la définition mais son contenu : la partition en
strates, l'existence de $i(x)$, et le fait que les $X_{i}$ sont les images
inverses des parties décroissantes. La section~4 ajoutera la moitié qui manque,
en montrant que $I' \mapsto X_{I'}$ commute aux réunions et aux intersections.}

Les mêmes définitions se propagent aux drapeaux :
\[
  \dot{X}_{\Delta_{0}} = \coprod_{i} \dot{X}_{i}, \qquad
  X^{*}_{\Delta_{0}} = X_{\Delta_{0}} \smallsetminus \dot{X}_{\Delta_{0}}
  = \coprod_{i} X_{i}^{*},
\]
et en degré $r$ par image inverse le long de
$\sigma_{r} \colon X_{\Delta_{r}} \to X_{\Delta_{0}}$, ce qui fait de
$\dot{X}_{\Delta_{r}}$ un fermé et de $X^{*}_{\Delta_{r}}$ un ouvert de
$X_{\Delta_{r}}$. Les applications $\alpha^{*}$ avec $\alpha(0) = 0$ respectent
la partition. En revanche — et l'avertissement est de l'auteur — les
$\dot{X}_{\Delta_{r}}$ ne forment pas un espace semi-simplicial : seule la moitié
des applications les préserve.

« Moralement », dit-il, $X^{*}_{\Delta_{r}}$ est la partie non singulière de
$X_{\Delta_{r}}$ et $\dot{X}_{\Delta_{r}}$ son bord. Le propos reste, dans
l'esprit de la section~1, de reconstituer $X$ à partir des strates $X_{i}^{*}$ —
la section~1 étant le cas où $I$ a deux éléments, $X_{0} = Y$ et $X_{1} = X$,
avec $Y = Y^{*}$ et $X^{*} = X \smallsetminus Y$, et où il avait fallu
introduire $\mathcal{V}_{Y,X}$.

\pagerange{56}{60}
\section*{3. Stratifications globales : introduction des tubes}

C'est le mot « il avait fallu » qui commande la section suivante. Pour tout
couple $i \leq j$ on pose
\[
  \overline{\mathcal{V}}_{i,j} = \mathcal{V}_{X_{i},X_{j}}
  \qquad (\text{voisinage tubulaire de } X_{i} \text{ dans } X_{j}),
\]
au sens laissé programmatique par la section~1, et l'on découpe dedans. Le
découpage se fait avec une seule partie fermée auxiliaire :
\[
  R_{i,j} = \bigcup \bigl\{ X_{k} \ \big|\ k < j,\ X_{k} \not\supset X_{i} \bigr\},
\]
la réunion des morceaux de $X_{j}$ qui ne contiennent pas $X_{i}$ — ceux qui
« ne voient pas » la strate depuis laquelle on
regarde.\footnote{Le manuscrit écrit cette partie de deux façons, $R^{*}_{i,j} =
\bigl(\bigcup_{X_{k} \not\supset X_{i}} X_{k}\bigr) \cap X_{j}$ et
$R_{i,j} = \bigcup_{k<j,\ X_{k} \cap X_{i} \neq X_{i}} X_{k}$, et emploie parfois
$R_{i}$ sans second indice pour la réunion non tronquée, qu'il note fermée dans
$X$. Les trois écritures désignent la même chose à la troncation près, et nous
n'en gardons qu'une.}

Le fait qui rend le découpage utilisable est le suivant, et il est démontré sur
la page :
\[
  R_{i,j} \cap X_{i} = \dot{X}_{i}.
\]
En effet, si $X_{k} \not\supset X_{i}$, l'intersection $X_{k} \cap X_{i}$ est,
par la condition d), la réunion des $X_{\ell}$ avec $\ell \leq k,i$ ; chacun de
ces $X_{\ell}$ est strictement contenu dans $X_{i}$, donc $\ell < i$ et
$X_{\ell} \subset \dot{X}_{i}$. L'inclusion inverse est triviale.

\subsection*{Trois espaces, et ce qu'ils valent}

On pose alors
\[
  \mathcal{V}_{i,j} = \overline{\mathcal{V}}_{i,j} \smallsetminus R_{i,j},
  \qquad
  \mathcal{V}^{*}_{i,j} = \overline{\mathcal{V}}_{i,j} \smallsetminus \dot{X}_{j}
  = \mathcal{V}_{i,j} \cap X_{j}^{*},
\]
la seconde égalité valant parce que $R_{i,j} \subset \dot{X}_{j}$. Et
l'identité ci-dessus donne
\[
  \mathcal{V}_{i,j} \cap X_{i} = X_{i}^{*} :
\]
$\mathcal{V}_{i,j}$ est un espace infinitésimal autour de la \emph{strate}
$X_{i}^{*}$, et non autour du fermé $X_{i}$. C'est le gain de tout le découpage.
Pour $i = j$ tout se dégénère proprement :
\[
  \overline{\mathcal{V}}_{i,i} = X_{i}, \qquad
  \mathcal{V}_{i,i} = \mathcal{V}^{*}_{i,i} = X_{i}^{*},
\]
de sorte que les strates elles-mêmes sont des cas particuliers des tubes.

Les pierres de construction sont donc les $\mathcal{V}_{i,j}$ et
$\mathcal{V}^{*}_{i,j}$, avec les inclusions incidentes :

\begin{tikzcd}[column sep=large, row sep=small]
& X_{i}^{*} \arrow[d, "\text{âme}"] \\
\mathcal{V}^{*}_{i,j} \arrow[r] \arrow[d] & \mathcal{V}_{i,j} \\
X_{j}^{*} &
\end{tikzcd}

$X_{i}^{*}$ étant l'\emph{âme} du tube : le mot est de l'auteur, et il est le
bon, puisque c'est la section sur laquelle le tube se rétracte.

\subsection*{La géométrie du tube}

Sous des hypothèses de lissité et de trivialité normale locale — pour une variété
algébrique complexe stratifiée par des sous-variétés fermées, par exemple — les
tubes prennent une forme précise, et c'est cette forme qui justifie qu'on les ait
préférés aux fermés. Le tube épointé $\mathcal{V}^{*}_{i,j}$ est un fibré
localement trivial sur $X_{i}^{*}$, de fibre un cylindre ouvert
$\mathbb{R} \times C_{i,j}$ où $C_{i,j}$ est une variété topologique compacte ;
et $\mathcal{V}_{i,j}$ est le fibré en cônes correspondant, de fibre le cône sur
$C_{i,j}$. Quand $j$ n'est pas un successeur de $i$, $C_{i,j}$ est encore une
variété, mais non compacte.\footnote{$C_{i,j}$ est ce qu'on appelle aujourd'hui
le \emph{lien} de la strate $X_{i}^{*}$ dans $X_{j}$, et la description « voisinage
$=$ fibré en cônes sur le lien » est la définition d'une stratification
\emph{conique}. La chose était disponible : les stratifications de
Thom (1969) et Mather (1970) ont exactement cette propriété, et Siebenmann (1972)
en a fait une définition. Le manuscrit ne cite personne et présente la propriété
comme une conséquence attendue de la trivialité normale locale.}

L'auteur juge alors ses pierres mal taillées : trop générales d'un côté, pas
assez de l'autre pour contenir les morceaux dont il aurait besoin, et il renvoie
à un autre de ses textes.\footnote{Le renvoi est à « Str$_{\infty}$, p.~84 »,
avec une formule $\mathcal{V}^{\,d'_{*}}_{d_{*},d''_{*}} \simeq
\widetilde{N}_{d_{*};(d'_{0},d''_{0})}|\,D^{*}_{d_{*}}$ dont aucun des symboles
n'est défini dans ce dossier. Nous n'identifions pas le texte visé.} La
restriction qu'il retient est de se borner aux couples $i \leq j$ où $j$ est un
\emph{successeur} de $i$ — c'est-à-dire où il n'existe aucun $k$ avec
$i < k < j$. Ce sont ces pierres-là, dites élémentaires, dont le lot suivant
fera le diagramme $\widetilde{I}$.

\subsection*{Les tubes mixtes}

Reste une famille intermédiaire, qui interpole entre $\mathcal{V}^{*}_{i,k}$ et
$\mathcal{V}_{i,k}$. Étant donné $i \leq j \leq k$, on pose
\[
  S_{j,k} = \bigcup \bigl\{ X_{\ell} \ \big|\ \ell < k,\ X_{\ell} \supset X_{j} \bigr\},
  \qquad
  \mathcal{V}^{\,j}_{i,k}
  = \overline{\mathcal{V}}_{i,k} \smallsetminus \bigl( R_{i,k} \cup S_{j,k} \bigr).
\]
Les deux parties $R_{j,k}$ et $S_{j,k}$ répondent à la partition de
$\{\ell \in I \mid \ell < k\}$ selon que $X_{\ell}$ contient ou non $X_{j}$, de
sorte que $R_{j,k} \cup S_{j,k} = \dot{X}_{k}$. Comme $i \leq j$ entraîne
$R_{i,k} \subset R_{j,k}$, on a bien $R_{i,k} \cup S_{j,k} \subset \dot{X}_{k}$,
et la famille est croissante en $j$ :
\[
  \mathcal{V}^{*}_{i,k} \;\subset\; \mathcal{V}^{\,j}_{i,k}
  \;\subset\; \mathcal{V}_{i,k},
  \qquad
  \mathcal{V}^{\,i}_{i,k} = \mathcal{V}^{*}_{i,k}, \qquad
  \mathcal{V}^{\,k}_{i,k} = \mathcal{V}_{i,k},
\]
les deux valeurs extrêmes venant de $R_{i,k} \cup S_{i,k} = \dot{X}_{k}$ d'une
part, de $S_{k,k} = \emptyset$ de l'autre. On retrouve pour $i = j = k$ la strate
$X_{i}^{*}$.\footnote{La page écrit $R_{j,k}$ là où nous écrivons $R_{i,k}$, et
c'est le seul endroit de la section où nous nous écartons d'elle. Avec
$R_{j,k}$, la réunion $R_{j,k} \cup S_{j,k}$ vaut $\dot{X}_{k}$ pour \emph{tout}
$j$, la famille $\mathcal{V}^{\,j}_{i,k}$ est constante et égale à
$\mathcal{V}^{*}_{i,k}$, et les deux identités que la page énonce juste en
dessous sont fausses. Avec $R_{i,k}$, elles sont vraies toutes les deux, la
famille interpole comme annoncé, et la note de marge — $S_{j,k} = \emptyset$ si
$j = k$, $S_{j,k} = X_{j}$ si $k$ est un successeur de $j$ — est exacte. La
correction est donc dictée par les énoncés que la page elle-même en tire. Le
manuscrit ajoute que $R_{j,k} \cap S_{j,k} = \dot{X}_{i}$ ; le membre de gauche
ne dépend pas de $i$, donc l'égalité ne peut se lire telle quelle, et nous ne
l'affirmons pas — ce que l'identité démontrée plus haut donne, c'est
$R_{j,k} \cap X_{j} = \dot{X}_{j}$.}

Le tube mixte $\mathcal{V}^{\,j}_{i,k}$ tend lui aussi à être un fibré localement
trivial de base $X_{i}^{*}$, mais sa fibre peut être singulière : dans le cas
$\mathcal{V}^{\,k}_{i,k} = \mathcal{V}_{i,k}$ avec $k$ successeur de $i$, la
singularité est isolée ; en général elle peut être de type arbitraire — il suffit
de prendre $X_{i} = X_{i}^{*}$ réduit à un point pour la mettre tout entière dans
la fibre. C'est ce qui diminue l'intérêt des tubes mixtes comme pierres, et
c'est pourquoi le programme annoncé à la fin de la section est celui-ci :
exprimer tout avec les $X_{i}^{*}$ et les tubes élémentaires
$\mathcal{V}_{i,j}$, $\mathcal{V}^{*}_{i,j}$ ($j$ successeur de $i$) ; voir
ensuite si certaines expressions se simplifient en admettant tous les
$i \leq j$ ; et seulement en dernier lieu, en vue du cas des diviseurs à
croisements normaux, voir comment les $\mathcal{V}^{\,j}_{i,k}$ s'insèrent dans
ce formalisme.

C'est ce programme que le dernier lot exécute, dans l'ordre annoncé, et qu'il
laisse inachevé à sa troisième étape.
\pagerange{61}{61}
\section*{Le diagramme fondamental}

On rassemble les tubes en une seule pièce par degré,
\[
  \mathcal{V}_{\Delta_{1}} = \coprod_{i \leq j} \mathcal{V}_{i,j},
  \qquad
  \mathcal{V}^{*}_{\Delta_{1}} = \coprod_{i \leq j} \mathcal{V}^{*}_{i,j}
  \ \subset \ \mathcal{V}_{\Delta_{1}} \ \text{(ouvert)},
\]
ce qui donne un diagramme global reflétant, pour chaque couple $i \leq j$, le
diagramme fondamental :

\begin{tikzcd}[column sep=large]
\mathcal{V}^{*}_{i,j} \arrow[r] \arrow[d] & \mathcal{V}_{i,j} \arrow[d, dashed] & X_{i}^{*} \arrow[l] \arrow[dl, dashed] \\
X_{j}^{*} & X &
\end{tikzcd}

Les flèches en pointillé sont celles que le recollement doit fournir. Et
l'auteur note aussitôt que le regroupement global n'est probablement pas utile,
car le carré qu'il forme n'a rien de cocartésien : le diagramme en trait plein
définit bien une limite inductive, mais elle n'a guère de chances d'être
homéomorphe à $X$ — dès que $X$ est connexe, par exemple, la somme disjointe sur
\emph{tous} les couples $i \leq j$ en compte trop.

L'observation est le ressort de tout ce qui suit : ce n'est pas un carré qu'il
faut, mais un diagramme dont la forme est celle de l'ensemble ordonné $I$ tout
entier, et il faut d'abord savoir de quelles flèches il dispose.

\pagerange{62}{63}
\section*{Quelles inclusions subsistent entre les pierres}

Les $\mathcal{V}^{\,j}_{i,k}$ sont des germes de sous-espaces de $X$, entre
lesquels les seuls morphismes canoniques sont des inclusions. La question est
donc de recenser lesquelles. Pour la traiter, deux hypothèses simplificatrices
sont posées :
\[
  X_{i}^{*} \neq \emptyset \quad \text{pour tout } i,
  \qquad X_{i} \subset X_{j} \implies i \leq j,
\]
la seconde faisant de $i \mapsto X_{i}$ un plongement d'ensembles
ordonnés ;\footnote{Ce sont exactement les deux hypothèses que la section 2
s'était refusée à faire, et la section 4 dira pourquoi elle avait raison : elles
ne survivent pas au passage aux images inverses. Elles ne servent ici qu'au
recensement, qui est une question locale de géométrie et non de fonctorialité.}
et une troisième, dite « raisonnable » :
$\overline{X_{i}^{*}} = X_{i}$ pour tout $i$ — la strate est dense dans son
fermé.

Le recensement donne alors ceci. Entre les strates elles-mêmes, aucune inclusion
pour des indices distincts. Entre tubes, $\mathcal{V}_{i,j} \subset
\mathcal{V}_{i',j'}$ force $i = i'$, puisqu'il faudrait $X_{i}^{*} \subset
X_{i'}^{*}$, puis $j = j'$ puisque le tube détermine la strate voisine dont il
est le voisinage. Entre tubes épointés, $\mathcal{V}^{*}_{i,j} \subset
\mathcal{V}^{*}_{i',j'}$ donne d'abord $i \leq i'$ par passage aux adhérences,
puis, comme $\mathcal{V}^{*}_{i,j} \subset X_{j}^{*}$ et
$\mathcal{V}^{*}_{i',j'} \subset X_{j'}^{*}$ sont non vides et se rencontrent,
$j = j'$ ; et enfin $i' = i$, puisque $i \leq i' < j$ et que $j$ est un
successeur de $i$.\footnote{Cette page est surchargée de reprises et la
transcription en donne le fil plutôt que la lettre ; plusieurs des mots qui
articulent l'argument y sont illisibles. La chaîne d'implications donnée ici est
celle que les deux petits diagrammes ordonnés de la page — $i$ dominé par $i'$ et
par $j$, $i'$ dominant $j'$ — permettent de suivre, et elle aboutit à la
conclusion que la page tire. Elle est donc une reconstruction, non une lecture.}

Restent les inclusions entre types différents, et il n'y en a que trois :
\[
  X_{i}^{*} \longrightarrow \mathcal{V}_{i,j}, \qquad
  \mathcal{V}^{*}_{i,j} \longrightarrow \mathcal{V}_{i,j}, \qquad
  \mathcal{V}^{*}_{i,j} \longrightarrow X_{j}^{*}.
\]
Ce sont exactement les trois flèches en trait plein du diagramme fondamental.
Rien d'autre n'est à envisager entre ces pièces élémentaires.\footnote{Les trois
inclusions portent au manuscrit des numéros que la transcription ne lit pas ;
nous les donnons sans numéro. Le fait que le recensement retombe précisément sur
le diagramme fondamental de la page 61 est ce qui autorise la construction
suivante : le diagramme $\widetilde{I}$ n'ajoute aucune flèche qui ne soit une
inclusion effective.}

\pagerange{64}{65}
\section*{Le diagramme $\widetilde{I}$}

À l'ensemble ordonné $I$ on associe un diagramme, noté $\widetilde{I}$, dont les
sommets sont des symboles
\[
  [i], \qquad [i,j], \qquad [i,j]^{*}
  \qquad (i,j \in I,\ j \text{ successeur de } i),
\]
et dont les flèches sont de trois types, toutes décrites par

\begin{tikzcd}[column sep=large, row sep=small]
{[i,j]}^{*} \arrow[r] \arrow[d] & {[i,j]} & {[i]}^{*} \arrow[l] \\
{[j]}^{*} & &
\end{tikzcd}

Notons $e$ le nombre de \emph{drapeaux élémentaires} — les couples $i < j$ avec
$j$ successeur de $i$. Alors $\widetilde{I}$ a $\mathrm{card}\,I + 2e$ sommets et
$3e$ flèches.\footnote{Le manuscrit note $\mathcal{V}$ l'ensemble des drapeaux
élémentaires, ce qui entre en collision avec les tubes $\mathcal{V}_{i,j}$ dans
la ligne même où les deux figurent. Nous écrivons $e$ pour le cardinal.} Si $I$
a deux éléments $i < j$, on retrouve le diagramme ci-dessus lui-même.

Trois exemples sont explicités, en remplaçant les symboles par les espaces
qu'ils désignent. Pour $I = \Delta_{n} = (0 < 1 < \cdots < n)$, une chaîne, on
trouve un escalier de $n$ marches :

\begin{tikzcd}[column sep=small, row sep=small, nodes={font=\scriptsize}]
X_{0}^{*} \arrow[r] & \mathcal{V}_{0,1} & \mathcal{V}^{*}_{0,1} \arrow[l] \arrow[d] & & & & \\
& & X_{1}^{*} \arrow[r] & \mathcal{V}_{1,2} & \mathcal{V}^{*}_{1,2} \arrow[l] \arrow[d] & & \\
& & & & X_{2}^{*} \arrow[r] & \mathcal{V}_{2,3} & \mathcal{V}^{*}_{2,3} \arrow[l] \arrow[d] \\
& & & & & & \cdots \arrow[d] \\
& & & & & & X_{n}^{*}
\end{tikzcd}

Pour $I$ formé d'un indice $0$ dominé par deux indices incomparables $a$ et $b$ :

\begin{tikzcd}[column sep=small, nodes={font=\scriptsize}]
\mathcal{V}^{*}_{0,b} \arrow[r] \arrow[d] & \mathcal{V}_{0,b} & X_{0}^{*} \arrow[l] \arrow[r] & \mathcal{V}_{0,a} & \mathcal{V}^{*}_{0,a} \arrow[l] \arrow[d] \\
X_{b}^{*} & & & & X_{a}^{*}
\end{tikzcd}

Et pour $I$ formé de deux indices $a$, $b$ incomparables dominés par un même $c$ :

\begin{tikzcd}[column sep=small, nodes={font=\scriptsize}]
X_{a}^{*} \arrow[d] & & X_{b}^{*} \arrow[d] \\
\mathcal{V}_{a,c} & & \mathcal{V}_{b,c} \\
\mathcal{V}^{*}_{a,c} \arrow[u] \arrow[r] & X_{c}^{*} & \mathcal{V}^{*}_{b,c} \arrow[u] \arrow[l]
\end{tikzcd}

Dans les trois cas les comptes se vérifient : sept sommets et six flèches pour
les deux derniers, $3n+1$ sommets et $3n$ flèches pour la chaîne.

La forme de ces diagrammes mérite d'être nommée. Chaque drapeau élémentaire
$i < j$ y contribue un motif $X_{i}^{*} \to \mathcal{V}_{i,j} \leftarrow
\mathcal{V}^{*}_{i,j} \to X_{j}^{*}$, c'est-à-dire un \emph{double cylindre
d'application} entre les deux strates, passant par le tube. C'est le motif
standard par lequel un espace stratifié se reconstruit à partir de ses strates et
de ses liens, et $\widetilde{I}$ n'est que l'assemblage de ces motifs le long de
l'ensemble ordonné.\footnote{En langue d'aujourd'hui : $\widetilde{I}$ est un
modèle $1$-catégorique de la catégorie des chemins sortants, et sa colimite
homotopique est ce qui calcule le type d'homotopie de l'espace stratifié. La
formulation générale, pour les stratifications coniques, est postérieure de trente
ans ; ce qui est ici est le diagramme, et l'assertion qu'il suffit.}

\pagerange{66}{67}
\section*{La fonction dimension}

Le diagramme se simplifie si l'on dispose sur $I$ d'une fonction à valeurs
entières
\[
  d \colon I \longrightarrow \mathbb{Z},
  \qquad \text{i.e.} \quad I = \coprod_{m} I_{m},
\]
jouant le rôle d'une dimension, et satisfaisant : pour $i \leq j$,
\[
  d(j) = d(i) + 1 \iff j \text{ est un successeur de } i.
\]
On peut alors regrouper par degré :
\[
  X_{m}^{*} = \coprod_{i \in I_{m}} X_{i}^{*},
  \qquad
  \mathcal{V}_{m,m+1} = \coprod_{\substack{i \leq j \\ d(i)=m,\ d(j)=m+1}}
    \mathcal{V}_{i,j},
  \qquad
  \mathcal{V}^{*}_{m,m+1} = \coprod_{\substack{i \leq j \\ d(i)=m,\ d(j)=m+1}}
    \mathcal{V}^{*}_{i,j},
\]
d'où un diagramme $(\mathrm{Diag})_{m}$ par degré,

\begin{tikzcd}[column sep=large, row sep=small]
X_{m}^{*} \arrow[r] & \mathcal{V}_{m,m+1} & \mathcal{V}^{*}_{m,m+1} \arrow[l] \arrow[d] \\
& & X^{*}_{m+1}
\end{tikzcd}

et l'escalier obtenu en les enchaînant, $I$ étant supposé fini et les sommets
occupés par des espaces vides n'étant pas écrits.

La variante par \emph{codimension} est la même à l'orientation près : on
considère alors les $\mathcal{V}_{d,d-1}$ et $\mathcal{V}^{*}_{d,d-1}$ au lieu
des $\mathcal{V}_{d,d+1}$ et $\mathcal{V}^{*}_{d,d+1}$, et l'escalier descend au
lieu de monter, de $X_{n}^{*}$ à $X_{0}^{*}$.\footnote{Le manuscrit caractérise
la fonction codimension par « $d(j) = j - 1$ si $j$ succ. de $i$ », ce qui
confond un indice et sa dimension. La condition voulue, parallèle à celle de la
fonction dimension, est $d(j) = d(i) - 1$ ; la transcription signale le lapsus
et ne le répare pas.}

\pagerange{68}{68}
\section*{Le théorème de recollement des tubes}

\textbf{Théorème} (première version). \emph{Considérons le diagramme de type
$\widetilde{I}$ formé par les $X_{i}^{*}$, les $\mathcal{V}^{*}_{i,j}$ et les
$\mathcal{V}_{i,j}$, pour $i,j \in I$ avec $j$ successeur de $i$, et par les
trois types de flèches en trait plein du diagramme fondamental. Alors $X$
s'identifie à la limite inductive de ce diagramme.}

Le mot que le manuscrit emploie pour l'opération est illisible ; « au topos »
est écrit puis barré juste avant, ce qui indique une hésitation sur la catégorie
où prendre la limite plutôt que sur la limite
elle-même.\footnote{Nous lisons « limite inductive » parce que c'est ce que toute
la section prépare, et parce que c'est ce que le cas de la section 1 donne. La
note marginale, également reprise plusieurs fois et en partie illisible, dit
qu'il faudrait préciser en quel sens $X$ est déterminé par $\mathrm{Top}(X)$ —
c'est-à-dire prendre la limite dans les topos plutôt que dans les espaces, ce
que le mot barré suggère aussi.}

L'énoncé se vérifie sur le cas dont tout le dossier est parti. Pour $I$ à deux
éléments $0 < 1$, avec $Y = X_{0}$ et $X = X_{1}$, on a $X_{0}^{*} = Y$,
$X_{1}^{*} = X \smallsetminus Y$, $\mathcal{V}_{0,1} = \mathcal{V}_{Y,X}$ et
$\mathcal{V}^{*}_{0,1} = \mathcal{V}^{*}_{Y,X}$ ; le diagramme $\widetilde{I}$
est alors
\[
  Y \longrightarrow \mathcal{V}_{Y,X} \longleftarrow \mathcal{V}^{*}_{Y,X}
  \longrightarrow X \smallsetminus Y,
\]
et sa limite inductive est exactement la somme amalgamée de la section 1, la
flèche $Y \to \mathcal{V}_{Y,X}$ n'y ajoutant rien puisqu'elle est une
équivalence d'homotopie. Le théorème général est donc la mise en famille de ce
cas-là, sur l'ensemble ordonné $I$ tout entier.

Aucune démonstration n'est donnée. L'auteur annonce qu'elle sera heuristique, et
qu'il veut d'abord un énoncé exprimant les « topos canoniques » associés à la
situation en termes d'un diagramme de type $\widetilde{I}$ formé de topos
élémentaires. C'est l'objet de la section suivante, et le dossier s'interrompt
avant d'y revenir.\footnote{Un paragraphe barré de deux longs traits diagonaux
précède cette annonce ; il proposait de démontrer le théorème en remplaçant les
voisinages tubulaires par de vrais voisinages ouverts, faute de disposer d'une
définition assez souple des premiers. La transcription en donne le squelette. Le
détour est significatif : c'est le même que la section 1 avait pris pour rendre
la somme amalgamée élémentaire.}

\pagerange{69}{71}
\section*{4. Topos canoniques : images inverses}

Les axiomes de la section 2 sont d'abord rappelés : les $X_{i}$ fermés et la
famille localement finie ; $i \leq j \Rightarrow X_{i} \subset X_{j}$ ;
$X = \bigcup X_{i}$ ; et $X_{i} \cap X_{j} = \bigcup_{k \leq i,j} X_{k}$.

\subsection*{Transport par image inverse}

Soit $X' \to X$ un espace au-dessus de $X$. La famille des
\[
  X'_{i} = X_{i} \times_{X} X'
\]
satisfait les mêmes conditions, et le système des drapeaux se transporte :
\[
  X'_{\Delta_{r}} \simeq X_{\Delta_{r}} \times_{X} X', \qquad
  X'^{\,*}_{\Delta_{r}} \simeq X^{*}_{\Delta_{r}} \times_{X} X', \qquad
  \dot{X}'_{\Delta_{r}} = \dot{X}_{\Delta_{r}} \times_{X} X'.
\]
Ces trois identités reposent sur le fait que le changement de base commute aux
réunions de sous-objets, donc à la définition de $\dot{X}_{i}$.

C'est ici que se paie une décision de la section 2. Après image inverse, un
$X'_{i}$ peut être vide, et deux indices distincts peuvent donner le même fermé :
c'est précisément pour pouvoir prendre ces images inverses qu'il n'était pas
commode de supposer les $X_{i}$ ou les $X_{i}^{*}$ non vides, ni que
$i \mapsto X_{i}$ soit un plongement d'ensembles ordonnés dans
$\mathcal{P}(X)$.\footnote{Le manuscrit le dit explicitement, page 70, et c'est
la justification rétrospective de la remarque b) de la section 2. On notera que
les hypothèses « raisonnables » du recensement des pages 62--63 sont exactement
celles-là : elles valent pour la stratification de départ, pas pour ses images
inverses.}

Pour les tubes, le transport demande davantage. Si $X' \to X$ est une immersion
ouverte, tout va de soi. Si c'est une immersion locale, on a encore
\[
  R'_{i,j} = R_{i,j} \times_{X} X', \qquad S'_{i,j} = S_{i,j} \times_{X} X',
\]
et, dans le cas propre, des isomorphismes
\[
  \mathcal{V}'_{i,j} \simeq \mathcal{V}_{i,j} \times_{X} X',
  \qquad
  \mathcal{V}'^{\,*}_{i,j} \simeq \mathcal{V}^{*}_{i,j} \times_{X} X',
  \qquad
  \mathcal{V}'^{\,j}_{i,k} \simeq \mathcal{V}^{\,j}_{i,k} \times_{X} X'.
\]
Le manuscrit est ici plus prudent qu'ailleurs, et il a raison de l'être : la
formation d'un voisinage tubulaire ne commute pas en général à la restriction à
une partie, et il faut que $X' \to X$ soit assez proche d'une immersion ouverte
pour que le tube d'une trace soit la trace du tube.\footnote{Le haut de la page
70 est traversé par une longue note marginale diagonale plusieurs fois reprise,
qui recouvre le texte ; la transcription en donne le fil. La restriction au cas
propre, portée en marge de la formule, est de l'auteur.}

\subsection*{Une conséquence : le théorème est local}

De ces changements de base il résulte que la démonstration du théorème de
recollement est une question \emph{locale} sur $X$. On peut donc se ramener à un
ouvert au-dessus duquel la famille est finie — la locale finitude est faite pour
cela — et supposer $I$ fini. C'est le seul progrès que le dossier fasse en
direction de la démonstration, et c'en est un vrai : il ramène un énoncé sur un
ensemble ordonné quelconque à un énoncé sur un ensemble ordonné fini.

\pagerange{71}{73}
\section*{Les cribles}

Une partie $I' \subset I$ est un \emph{crible} si elle est close vers le bas :
\[
  i \leq j \in I' \implies i \in I'.
\]
On lui associe la partie fermée
\[
  X_{I'} = \bigcup_{i \in I'} X_{i}.
\]
La correspondance est monotone, et elle respecte les deux opérations :
\[
  X_{\bigcup_{\alpha} I'_{\alpha}} = \bigcup_{\alpha} X_{I'_{\alpha}},
  \qquad X_{\emptyset} = \emptyset,
  \qquad
  X_{I' \cap I''} = X_{I'} \cap X_{I''}, \qquad X_{I} = X,
\]
la troisième venant de la condition d) des axiomes — et de nulle part
ailleurs.\footnote{C'est la moitié qui manquait à la section 2. Les cribles de
$I$ sont les fermés de la topologie d'Alexandrov sur $I$, celle dont les ouverts
sont les parties croissantes ; dire que $I' \mapsto X_{I'}$ commute aux réunions
quelconques et aux intersections finies, c'est dire que $x \mapsto i(x)$ est une
application continue $X \to I$ pour cette topologie. La condition d), qui
paraissait technique, est exactement ce qui rend la stratification continue.}

Prenant $X' = X_{I'}$, on obtient sur $X'$ une stratification de type $I'$, et
les parties-cribles de $X'$ pour cette stratification sont les traces sur $X'$
des parties-cribles de $X$. Les espaces élémentaires y sont les $X_{i}^{*}$, les
$\mathcal{V}^{*}_{i,j}$ et les $\mathcal{V}_{i,j}$ pour $i,j \in I'$.

Un cas particulier vaut d'être isolé, parce qu'il justifie une désinvolture
antérieure. Soit $I_{0}$ l'ensemble des $i$ tels que $X_{i} = \emptyset$ : c'est
un crible, et tous les espaces qui lui sont associés sont vides. Le diagramme
$\widetilde{I}$, une fois les sommets vides négligés, ne change donc pas quand on
remplace $I$ par $I \smallsetminus I'_{0}$ pour un crible $I'_{0} \subset I_{0}$.
Plus généralement, pour tout crible $I'$ de $I$, le diagramme
$\widetilde{I \smallsetminus I'}$ est contenu dans $\widetilde{I}$ : le
complémentaire d'un crible est clos vers le haut, donc un couple qui y est
élémentaire l'était déjà dans $I$.

\subsection*{Ouverts entre deux cribles}

Soient maintenant deux cribles emboîtés $I'' \subset I' \subset I$, d'où
$X_{I''} \subset X_{I'}$ et la partie localement fermée
\[
  \mathcal{U}_{I',I''} = X_{I'} \smallsetminus X_{I''}.
\]
On peut y regarder la stratification de type $I' \smallsetminus I''$, définie par
\[
  X_{i',\,\mathcal{U}_{I',I''}}
  = X_{i'} \smallsetminus \bigl(X_{i'} \cap X_{I''}\bigr)
  \ \subset \ \mathcal{U}_{I',I''},
\]
dont les espaces élémentaires sont les $X_{i'}^{*}$ pour
$i' \in I' \smallsetminus I''$ et les tubes associés aux couples $(i',j')$ avec
$j'$ successeur de $i'$. Ceux-ci se comparent aux tubes de $X$ par un carré
d'inclusions,

\begin{tikzcd}[column sep=large, row sep=small]
\mathcal{V}_{i',j',X'} \arrow[r, hook] & \mathcal{V}_{i',j'} \\
\mathcal{V}^{*}_{i',j',X'} \arrow[r, hook] \arrow[u, hook] & \mathcal{V}^{*}_{i',j'} \arrow[u, hook]
\end{tikzcd}

mais — et l'avertissement est de l'auteur, il vaut d'être conservé tel quel —
il n'est pas clair en général que ces inclusions soient des équivalences
d'homotopie. Mieux vaut donc ne pas identifier trop vite les constructions faites
sur un $X_{I'}$ et celles faites sur un $\mathcal{U}_{I',I''}$.

\pagerange{74}{74}
\section*{Où le dossier s'arrête}

La partie C de la section 4 s'ouvre sur la promesse d'une famille
$\mathcal{V}^{\,I'''}_{I',I''}$, indexée par trois cribles
\[
  I',\ I''' \subset I'' \subset I,
  \qquad\text{d'où}\qquad
  X_{I'},\ X_{I'''} \subset X_{I''} \hookrightarrow X_{I} = X,
\]
et le dossier s'arrête ici. La page 74 est la dernière du fonds pour cette cote,
et rien n'indique que la suite existe ailleurs.

La construction est donc inachevée, et cette édition ne lui invente pas de fin.
Ce qui se dit sans risque tient à la notation seule : les trois cribles occupent
les places qu'occupaient les trois indices $i \leq j \leq k$ dans les tubes
mixtes $\mathcal{V}^{\,j}_{i,k}$ de la section 3, l'exposant portant l'indice
médian ; et le programme énoncé à la fin de cette même section plaçait les tubes
mixtes en dernier, « eu égard aux cas des diviseurs à croisements normaux ».
C'est cette dernière étape que la page 74 commence. Ce qu'elle aurait donné, on
ne le lit pas ici.

\end{document}
